% Part 1
\chapter{Part 1: Personal \& General Topics}

\section*{PART 1: PERSONAL \& LIFESTYLE TOPICS}
\addcontentsline{toc}{section}{PART 1: PERSONAL \& LIFESTYLE TOPICS}

\section{Work or Studies}
\begin{ieltsquestion}
\textbf{What work do you do?}
\end{ieltsquestion}
\textbf{Sample 1: Student Persona (Band 7-8)}

At the moment, I am not working full-time because I am a second-year university student. Most of my weekdays are taken up by lectures, assignments, and a few club activities, so studying is really my main focus for now.

\textbf{Sample 2: Working Professional Persona (Band 7-8)}

I work as a marketing executive for a mid-sized education company. My job mainly involves planning social media campaigns, writing content, and tracking how well each campaign performs.

\separator

\begin{ieltsquestion}
\textbf{What subjects are you studying?}
\end{ieltsquestion}
\textbf{Sample 1: Student Persona (Band 7-8)}

I am majoring in business administration, so I study subjects like marketing, finance, management, and consumer behaviour. I would say marketing is the most enjoyable one because it allows me to be both analytical and creative.

\textbf{Sample 2: Working Professional Persona (Band 7-8)}

I am not studying at university anymore, but I still take short online courses related to digital marketing and data analysis. They help me stay up to date with industry trends, which is pretty important in my field.

\separator

\begin{ieltsquestion}
\textbf{Do you like your subject?}
\end{ieltsquestion}
\textbf{Sample 1: Student Persona (Band 7-8)}

Yes, I do, although it can be quite demanding at times. What I like most is that the theories are not too abstract; I can actually apply them to real businesses and even small projects with my classmates.

\textbf{Sample 2: Working Professional Persona (Band 7-8)}

Looking back, I did enjoy my major because it gave me a solid foundation for my current job. Some subjects were dry, especially statistics, but they turned out to be more useful than I expected.

\separator

\begin{ieltsquestion}
\textbf{Do you like your job?}
\end{ieltsquestion}
\textbf{Sample 1: Student Persona (Band 7-8)}

Since I am still a student, I do not have a full-time job yet. I do some freelance design work occasionally, and I quite enjoy it because it gives me a taste of professional life without overwhelming my study schedule.

\textbf{Sample 2: Working Professional Persona (Band 7-8)}

Yes, for the most part. The workload can be intense when deadlines pile up, but I like the fast-paced environment and the feeling of seeing an idea turn into a real campaign.

\separator

\begin{ieltsquestion}
\textbf{Why did you choose to study that subject?}
\end{ieltsquestion}
\textbf{Sample 1: Student Persona (Band 7-8)}

I chose business administration because I have always been curious about how companies attract customers and make decisions. Also, it is a fairly versatile major, so it gives me several career options instead of locking me into one path.

\textbf{Sample 2: Working Professional Persona (Band 7-8)}

I chose marketing-related subjects because I was interested in both creativity and psychology. I liked the idea of understanding what people want and then turning that insight into a message that actually connects with them.

\separator

\begin{ieltsquestion}
\textbf{Why did you choose to do that type of work (or that job)?}
\end{ieltsquestion}
\textbf{Sample 1: Student Persona (Band 7-8)}

For my freelance work, I chose graphic design because it fits naturally with my interests. It is flexible, and I can take on small projects during quieter weeks without affecting my classes too much.

\textbf{Sample 2: Working Professional Persona (Band 7-8)}

I chose marketing because it is a field where things never stay the same for long. I enjoy that mix of strategy, creativity, and measurable results, so the job keeps me on my toes in a good way.

\separator

\begin{ieltsquestion}
\textbf{What requirements did you need to meet to get your current job?}
\end{ieltsquestion}
\textbf{Sample 1: Student Persona (Band 7-8)}

For the small freelance projects I do, I mainly needed a decent portfolio and the ability to communicate clearly with clients. Grades were not the main issue; people cared more about whether I could deliver clean designs on time.

\textbf{Sample 2: Working Professional Persona (Band 7-8)}

I needed a bachelor's degree, some internship experience, and a basic understanding of digital tools like Google Analytics and Meta Ads. During the interview, they also cared a lot about problem-solving skills and whether I could work under pressure.

\separator

\begin{ieltsquestion}
\textbf{Do you think that your subject is popular in your country?}
\end{ieltsquestion}
\textbf{Sample 1: Student Persona (Band 7-8)}

Yes, business-related majors are very popular in Vietnam. A lot of students choose them because they seem practical and open the door to different careers, from banking to marketing to entrepreneurship.

\textbf{Sample 2: Working Professional Persona (Band 7-8)}

Definitely. Marketing and business courses are in high demand because many companies are moving online and need people who understand customers, branding, and digital platforms.

\separator

\begin{ieltsquestion}
\textbf{Do you have any plans for your studies in the next five years?}
\end{ieltsquestion}
\textbf{Sample 1: Student Persona (Band 7-8)}

Yes, my immediate plan is to finish my degree with a strong GPA and complete at least one internship before graduation. After that, I may pursue a master's degree abroad if I can get a scholarship.

\textbf{Sample 2: Working Professional Persona (Band 7-8)}

I am not planning to go back to university full-time, but I do want to keep learning through professional certificates. In the next few years, I would like to improve my data skills because marketing is becoming much more numbers-driven.

\separator

\begin{ieltsquestion}
\textbf{Do you have any plans for your work in the next five years?}
\end{ieltsquestion}
\textbf{Sample 1: Student Persona (Band 7-8)}

I hope to get an entry-level job in marketing after I graduate. Ideally, I would like to work for a company where I can learn from experienced colleagues and gradually build a clear career path.

\textbf{Sample 2: Working Professional Persona (Band 7-8)}

Yes, I would like to move into a senior role where I can manage a small team and handle bigger campaigns. I am not in a rush, but I want to keep building both technical skills and leadership experience.

\separator

\begin{ieltsquestion}
\textbf{Do you prefer to study in the mornings or in the afternoons?}
\end{ieltsquestion}
\textbf{Sample 1: Student Persona (Band 7-8)}

I prefer studying in the morning because my mind feels fresher then. If I leave difficult tasks until the afternoon, I tend to lose focus and get distracted by messages or other plans.

\textbf{Sample 2: Working Professional Persona (Band 7-8)}

When I take online courses, I prefer studying in the evening after work. It is not always ideal because I can be tired, but it is the only time when I can sit down without work meetings interrupting me.

\separator

\begin{ieltsquestion}
\textbf{How much time do you spend on your studies each week?}
\end{ieltsquestion}
\textbf{Sample 1: Student Persona (Band 7-8)}

I probably spend around twenty to twenty-five hours a week outside class reviewing lessons and doing assignments. The number goes up sharply before exams because I have to revise, prepare presentations, and finish group projects.

\textbf{Sample 2: Working Professional Persona (Band 7-8)}

For professional learning, I usually spend three or four hours a week. It is not a huge amount, but I try to be consistent because small, regular learning sessions are easier to maintain than occasional cramming.

\separator

\begin{ieltsquestion}
\textbf{Do you want to change your major?}
\end{ieltsquestion}
\textbf{Sample 1: Student Persona (Band 7-8)}

No, not really. There are difficult moments, of course, but overall I still feel that my major suits my strengths and gives me enough room to explore different interests.

\textbf{Sample 2: Working Professional Persona (Band 7-8)}

If I could go back, I might add more technology-related subjects, but I would not completely change my major. Business gave me a useful base, and I can always fill the gaps through short courses.

\separator

\begin{ieltsquestion}
\textbf{Do you want to change to another job?}
\end{ieltsquestion}
\textbf{Sample 1: Student Persona (Band 7-8)}

At this stage, I am not tied to one job yet, so I am still exploring. I would like to try internships in a few different areas before deciding what kind of work really suits me.

\textbf{Sample 2: Working Professional Persona (Band 7-8)}

Not right now. I still feel I am learning a lot in my current position, but in the long run I might move to a larger company if it offers better training and career progression.

\separator

\begin{ieltsquestion}
\textbf{Are you looking forward to working?}
\end{ieltsquestion}
\textbf{Sample 1: Student Persona (Band 7-8)}

Yes, I am, although I know it will be challenging. I am looking forward to becoming more independent, earning my own income, and applying what I have learned in a real workplace.

\textbf{Sample 2: Working Professional Persona (Band 7-8)}

I already work, but I still look forward to new projects because they keep the job interesting. Whenever I get to launch something fresh, it gives me a sense of momentum.

\separator

\begin{ieltsquestion}
\textbf{Do you miss being a student?}
\end{ieltsquestion}
\textbf{Sample 1: Student Persona (Band 7-8)}

I am still a student, so maybe I do not appreciate it fully yet. But I think I will miss the flexibility, the long breaks, and the chance to spend so much time with friends.

\textbf{Sample 2: Working Professional Persona (Band 7-8)}

Yes, sometimes. Student life was stressful in its own way, but I miss having fewer responsibilities and being able to learn without worrying too much about bills or workplace pressure.

\separator

\begin{ieltsquestion}
\textbf{What technology do you use when you study?}
\end{ieltsquestion}
\textbf{Sample 1: Student Persona (Band 7-8)}

I use my laptop for almost everything, from taking notes to making slides and doing research. I also rely on apps like Google Drive and Notion because they help me organise materials and collaborate with classmates.

\textbf{Sample 2: Working Professional Persona (Band 7-8)}

When I study after work, I mostly use online learning platforms, note-taking apps, and sometimes AI tools to summarise long articles. They save me a lot of time, especially when I am already tired after a full day.

\separator

\begin{ieltsquestion}
\textbf{What technology do you use at work?}
\end{ieltsquestion}
\textbf{Sample 1: Student Persona (Band 7-8)}

For my freelance design tasks, I mainly use design software like Canva and Adobe Illustrator. I also use messaging apps to communicate with clients, which makes the whole process much smoother.

\textbf{Sample 2: Working Professional Persona (Band 7-8)}

I use quite a few digital tools at work, including project management software, analytics dashboards, and content scheduling platforms. Without them, it would be much harder to coordinate tasks and track campaign performance.

\separator

\begin{ieltsquestion}
\textbf{What changes would you like to see in your school?}
\end{ieltsquestion}
\textbf{Sample 1: Student Persona (Band 7-8)}

I would like my university to offer more practical workshops and company visits. The lectures are useful, but students also need more exposure to real business situations before they graduate.

\textbf{Sample 2: Working Professional Persona (Band 7-8)}

Thinking back to my university, I wish it had provided stronger career guidance. Many students, including me at the time, knew the theory but felt unsure about how to turn it into an actual career plan.

\separator

\begin{ieltsquestion}
\textbf{Who helps you the most? And how?}
\end{ieltsquestion}
\textbf{Sample 1: Student Persona (Band 7-8)}

My closest friend helps me the most because we study the same major and often prepare for exams together. She explains difficult concepts in a simple way, and she also keeps me motivated when I feel overwhelmed.

\textbf{Sample 2: Working Professional Persona (Band 7-8)}

At work, my line manager helps me the most. She gives clear feedback on my campaigns, points out blind spots, and encourages me to take more ownership instead of just following instructions.

\separator

\begin{ieltsquestion}
\textbf{What do you think is the most important at the moment?}
\end{ieltsquestion}
\textbf{Sample 1: Student Persona (Band 7-8)}

Right now, the most important thing for me is building a solid foundation. I want to keep my grades stable, improve my English, and gain some practical experience before entering the job market.

\textbf{Sample 2: Working Professional Persona (Band 7-8)}

For me, the priority is maintaining a healthy balance between performance and well-being. I want to do well at work, but I also know that burning out would hurt my career in the long run.

\separator

\begin{ieltsquestion}
\textbf{What are the benefits of being your age?}
\end{ieltsquestion}
\textbf{Sample 1: Student Persona (Band 7-8)}

I think the biggest benefit is having time to experiment. I am young enough to make mistakes, try different activities, and change direction without feeling that it is too late.

\textbf{Sample 2: Working Professional Persona (Band 7-8)}

At my age, I have more independence than I did as a student, but I am still young enough to learn quickly and take risks. It feels like a good stage for building habits and making long-term career decisions.

\separator


\vspace{0.8em}
\begin{center}
\renewcommand{\arraystretch}{1.65}
\rowcolors{2}{ForumFreshLight}{white}
\begin{tabularx}{\textwidth}{|V{0.22\textwidth}|C{0.07\textwidth}|I{0.20\textwidth}|Y|V{0.18\textwidth}|}
\hline
\rowcolor{ForumFreshBlue!20}
\textbf{Từ} & \textbf{Loại} & \textbf{IPA} & \textbf{English Meaning} & \textbf{Dịch nghĩa} \\
\hline
\textbf{take up} & phr v & /teɪk ʌp/ & to use or occupy time or space & chiếm thời gian/không gian \\ \hline
\textbf{solid foundation} & n phr & /ˈsɒl.ɪd faʊnˈdeɪ.ʃən/ & a strong basic understanding & nền tảng vững chắc \\ \hline
\textbf{versatile} & adj & /ˈvɜː.sə.taɪl/ & useful in many different ways & linh hoạt, đa năng \\ \hline
\textbf{keep someone on their toes} & idiom & /kiːp ˈsʌm.wʌn ɒn ðeə təʊz/ & to make someone stay alert and active & khiến ai luôn tập trung, nhanh nhạy \\ \hline
\textbf{portfolio} & n & /pɔːtˈfəʊ.li.əʊ/ & a collection of work samples & hồ sơ sản phẩm/năng lực \\ \hline
\textbf{entry-level} & adj & /ˈen.tri ˌlev.əl/ & suitable for someone starting a career & ở cấp độ mới vào nghề \\ \hline
\textbf{career progression} & n phr & /kəˈrɪə prəˈɡreʃ.ən/ & gradual advancement in a career & sự thăng tiến nghề nghiệp \\ \hline
\textbf{collaborate} & v & /kəˈlæb.ə.reɪt/ & to work together with others & cộng tác \\ \hline
\textbf{career guidance} & n phr & /kəˈrɪə ˈɡaɪ.dəns/ & advice about jobs and career paths & định hướng nghề nghiệp \\ \hline
\textbf{burn out} & phr v & /bɜːn aʊt/ & to become exhausted from too much work & kiệt sức vì làm việc quá nhiều \\ \hline
\end{tabularx}
\end{center}

\newpage

\section{Hometown}
\begin{ieltsquestion}
\textbf{Where is your hometown?}
\end{ieltsquestion}
\textbf{Sample 1: Student Persona (Band 7-8)}

My hometown is Da Nang, a coastal city in central Vietnam. My family still lives there, and I usually go back during semester breaks, so it still feels very much like home to me.

\textbf{Sample 2: Working Professional Persona (Band 7-8)}

I come from Da Nang, which is one of the most liveable cities in Vietnam. Although I now work in Ho Chi Minh City, I still consider Da Nang my real hometown because most of my family and childhood memories are there.

\separator

\begin{ieltsquestion}
\textbf{Is that a big city or a small place?}
\end{ieltsquestion}
\textbf{Sample 1: Student Persona (Band 7-8)}

I would say it is a medium-sized city. It is not as hectic as Hanoi or Ho Chi Minh City, but it is definitely not a sleepy town either, especially during the tourist season.

\textbf{Sample 2: Working Professional Persona (Band 7-8)}

It is somewhere in between. Da Nang has modern infrastructure and a growing economy, but it still feels more manageable and less crowded than the major business hubs in Vietnam.

\separator

\begin{ieltsquestion}
\textbf{Please describe your hometown a little.}
\end{ieltsquestion}
\textbf{Sample 1: Student Persona (Band 7-8)}

Da Nang is known for its beaches, bridges, and relaxed lifestyle. What I like is that you can go from a busy city street to the seaside in just a few minutes, which makes daily life feel quite balanced.

\textbf{Sample 2: Working Professional Persona (Band 7-8)}

It is a clean and well-planned city with a nice mix of nature and urban development. There are beaches, mountains, good restaurants, and enough business activity to make it feel dynamic without being overwhelming.

\separator

\begin{ieltsquestion}
\textbf{Do you like your hometown?}
\end{ieltsquestion}
\textbf{Sample 1: Student Persona (Band 7-8)}

Yes, absolutely. I like it because it is peaceful enough for studying and family life, but it still has plenty of cafes, shopping areas, and places to hang out with friends.

\textbf{Sample 2: Working Professional Persona (Band 7-8)}

Yes, I do. After living in a bigger city for work, I appreciate Da Nang even more because the pace of life is calmer and people seem less stressed on a day-to-day basis.

\separator

\begin{ieltsquestion}
\textbf{What do you like (most) about your hometown?}
\end{ieltsquestion}
\textbf{Sample 1: Student Persona (Band 7-8)}

What I like most is the beach. Whenever I feel stressed about exams or assignments, I can go there for a walk, get some fresh air, and clear my head quite quickly.

\textbf{Sample 2: Working Professional Persona (Band 7-8)}

For me, the best thing is the quality of life. The city has good food, beautiful scenery, and shorter commuting times, so it is much easier to maintain a healthy work-life balance there.

\separator

\begin{ieltsquestion}
\textbf{Is there anything you dislike about it?}
\end{ieltsquestion}
\textbf{Sample 1: Student Persona (Band 7-8)}

If I had to choose one thing, I would say the weather can be a bit harsh. During the rainy season, heavy storms and flooding sometimes disrupt classes and make travelling around the city inconvenient.

\textbf{Sample 2: Working Professional Persona (Band 7-8)}

The main drawback is that there are still fewer career opportunities compared with bigger cities. For some specialised fields, people often have to move away if they want faster career growth.

\separator

\begin{ieltsquestion}
\textbf{How long have you been living there?}
\end{ieltsquestion}
\textbf{Sample 1: Student Persona (Band 7-8)}

I lived there for most of my childhood before moving away for university. So even though I am not there full-time now, I spent around eighteen years growing up in that city.

\textbf{Sample 2: Working Professional Persona (Band 7-8)}

I lived there until I finished university, so roughly twenty-two years. I moved to Ho Chi Minh City for work a few years ago, but I still visit Da Nang whenever I get a long weekend.

\separator

\begin{ieltsquestion}
\textbf{Do you like living there?}
\end{ieltsquestion}
\textbf{Sample 1: Student Persona (Band 7-8)}

Yes, I really do. Life there is comfortable, and I like the fact that most places are not too far from each other, so I can get around easily without wasting too much time.

\textbf{Sample 2: Working Professional Persona (Band 7-8)}

Yes, living there is quite pleasant. If my job allowed remote work permanently, I would seriously consider moving back because the city offers a calmer lifestyle and lower living costs.

\separator

\begin{ieltsquestion}
\textbf{Do you think you will continue living there for a long time?}
\end{ieltsquestion}
\textbf{Sample 1: Student Persona (Band 7-8)}

I am not completely sure. I may need to move to a bigger city after graduation for better job opportunities, but in the long run, I would love to return to Da Nang and settle down there.

\textbf{Sample 2: Working Professional Persona (Band 7-8)}

Probably not in the immediate future because my current job is based elsewhere. However, I can imagine moving back later in life, especially if I can find a suitable role or start my own business there.

\separator

\begin{ieltsquestion}
\textbf{What's your hometown famous for?}
\end{ieltsquestion}
\textbf{Sample 1: Student Persona (Band 7-8)}

It is famous for My Khe Beach, the Dragon Bridge, and local food like Mi Quang. Tourists also use Da Nang as a gateway to Hoi An and Ba Na Hills, so it is quite well-known now.

\textbf{Sample 2: Working Professional Persona (Band 7-8)}

Da Nang is famous for being clean, scenic, and tourist-friendly. In recent years, it has also built a reputation as a good place for conferences, tourism services, and small start-ups.

\separator

\begin{ieltsquestion}
\textbf{Did you learn about the culture of your hometown in your childhood?}
\end{ieltsquestion}
\textbf{Sample 1: Student Persona (Band 7-8)}

Yes, mostly through family activities and school events. For example, I learned about local festivals, traditional food, and the way people in central Vietnam value resilience and hospitality.

\textbf{Sample 2: Working Professional Persona (Band 7-8)}

Yes, I did, but in a very natural way rather than through formal lessons. Growing up there, I absorbed the local culture through food, family gatherings, festivals, and everyday conversations with older relatives.

\separator

\begin{ieltsquestion}
\textbf{Did you learn about the history of your hometown at school?}
\end{ieltsquestion}
\textbf{Sample 1: Student Persona (Band 7-8)}

Yes, but only to some extent. We learned about Da Nang's role as a port city and some historical sites nearby, but I think the lessons were quite brief and textbook-based.

\textbf{Sample 2: Working Professional Persona (Band 7-8)}

Yes, I remember learning about its development from a coastal trading area into a modern city. To be honest, I did not find it very exciting as a student, but I appreciate it more now.

\separator

\begin{ieltsquestion}
\textbf{Have you learned anything about the history of your hometown?}
\end{ieltsquestion}
\textbf{Sample 1: Student Persona (Band 7-8)}

Recently, I learned more about how quickly Da Nang has changed over the past few decades. It used to be much quieter, but now it has become a major tourist and transport centre.

\textbf{Sample 2: Working Professional Persona (Band 7-8)}

Yes, especially after I started paying attention to urban development. I learned that the city's growth was not accidental; it came from long-term investment in infrastructure and tourism.

\separator

\begin{ieltsquestion}
\textbf{Are there many young people in your hometown?}
\end{ieltsquestion}
\textbf{Sample 1: Student Persona (Band 7-8)}

Yes, there are quite a lot of young people, especially around universities, cafes, and entertainment areas. However, some still move to bigger cities after graduation because they want broader career options.

\textbf{Sample 2: Working Professional Persona (Band 7-8)}

There are many young people, but the city also loses some talent to Hanoi and Ho Chi Minh City. That said, more young professionals are staying now because tourism, technology, and service industries are growing.

\separator

\begin{ieltsquestion}
\textbf{Is your hometown a good place for young people to pursue their careers?}
\end{ieltsquestion}
\textbf{Sample 1: Student Persona (Band 7-8)}

It depends on the field. For tourism, hospitality, and some creative jobs, Da Nang can be a promising place, but for highly specialised careers, young people may still need to look elsewhere.

\textbf{Sample 2: Working Professional Persona (Band 7-8)}

I think it is becoming a better place for young professionals. The market is not as large as in Ho Chi Minh City, but the lower living costs and improving business environment make it increasingly attractive.

\separator


\vspace{0.8em}
\begin{center}
\renewcommand{\arraystretch}{1.65}
\rowcolors{2}{ForumFreshLight}{white}
\begin{tabularx}{\textwidth}{|V{0.22\textwidth}|C{0.07\textwidth}|I{0.20\textwidth}|Y|V{0.18\textwidth}|}
\hline
\rowcolor{ForumFreshBlue!20}
\textbf{Từ} & \textbf{Loại} & \textbf{IPA} & \textbf{English Meaning} & \textbf{Dịch nghĩa} \\
\hline
\textbf{coastal city} & n phr & /ˈkəʊ.stəl ˈsɪt.i/ & a city located near the sea & thành phố ven biển \\ \hline
\textbf{liveable} & adj & /ˈlɪv.ə.bəl/ & pleasant and suitable to live in & đáng sống \\ \hline
\textbf{hectic} & adj & /ˈhek.tɪk/ & very busy and full of activity & bận rộn, hối hả \\ \hline
\textbf{manageable} & adj & /ˈmæn.ɪ.dʒə.bəl/ & easy enough to control or deal with & dễ quản lý, vừa tầm \\ \hline
\textbf{work-life balance} & n phr & /ˌwɜːk laɪf ˈbæl.əns/ & a healthy division between work and personal life & cân bằng công việc và cuộc sống \\ \hline
\textbf{drawback} & n & /ˈdrɔː.bæk/ & a disadvantage or problem & điểm bất lợi \\ \hline
\textbf{settle down} & phr v & /ˈset.əl daʊn/ & to begin living a stable life in one place & ổn định cuộc sống \\ \hline
\textbf{gateway to} & n phr & /ˈɡeɪt.weɪ tuː/ & a place that gives access to another place & cửa ngõ dẫn đến \\ \hline
\textbf{resilience} & n & /rɪˈzɪl.i.əns/ & the ability to recover from difficulties & sự kiên cường \\ \hline
\textbf{urban development} & n phr & /ˈɜː.bən dɪˈvel.əp.mənt/ & the growth and planning of a city & phát triển đô thị \\ \hline
\end{tabularx}
\end{center}

\newpage

\section{Home and Accommodation}
\begin{ieltsquestion}
\textbf{Who do you live with?}
\end{ieltsquestion}
\textbf{Sample 1: Student Persona (Band 7-8)}

At the moment, I live with two classmates in a small rented apartment near our university. It is not luxurious, but sharing the place helps us cut down on living costs.

\textbf{Sample 2: Working Professional Persona (Band 7-8)}

I live by myself in a rented apartment, although my family visits me from time to time. Living alone gives me more privacy, which I really appreciate after a long day at work.

\separator

\begin{ieltsquestion}
\textbf{Do you live in an apartment or a house?}
\end{ieltsquestion}
\textbf{Sample 1: Student Persona (Band 7-8)}

I live in an apartment. It is quite compact, but it has everything I need as a student, like a desk, a small kitchen, and a decent internet connection.

\textbf{Sample 2: Working Professional Persona (Band 7-8)}

I live in an apartment building close to my office. It is convenient because there is security, parking, and a small grocery store downstairs.

\separator

\begin{ieltsquestion}
\textbf{Do you prefer living in a house or an apartment?}
\end{ieltsquestion}
\textbf{Sample 1: Student Persona (Band 7-8)}

For now, I prefer an apartment because it is easier to maintain and usually closer to campus. In the future, though, I might prefer a house if I have a family.

\textbf{Sample 2: Working Professional Persona (Band 7-8)}

At this stage of my life, I prefer an apartment. It suits my busy schedule because I do not have to spend much time on repairs, cleaning, or looking after a garden.

\separator

\begin{ieltsquestion}
\textbf{What kinds of accommodation do you live in?}
\end{ieltsquestion}
\textbf{Sample 1: Student Persona (Band 7-8)}

I live in a shared student apartment. Each of us has a small bedroom, and we share the kitchen and living area, so it feels both affordable and sociable.

\textbf{Sample 2: Working Professional Persona (Band 7-8)}

I live in a one-bedroom serviced apartment. It is fairly modern and low-maintenance, which is exactly what I need because my work schedule can be unpredictable.

\separator

\begin{ieltsquestion}
\textbf{Can you describe the place where you live?}
\end{ieltsquestion}
\textbf{Sample 1: Student Persona (Band 7-8)}

It is a modest apartment on the fourth floor of an older building. The furniture is simple, but the place gets plenty of natural light, so it feels more comfortable than it looks at first.

\textbf{Sample 2: Working Professional Persona (Band 7-8)}

My apartment is small but well-organised. It has a bedroom, a living room, and a little balcony where I keep a few plants, which makes the place feel less boxed in.

\separator

\begin{ieltsquestion}
\textbf{Please describe the room you live in.}
\end{ieltsquestion}
\textbf{Sample 1: Student Persona (Band 7-8)}

My room is quite small, but I have tried to make it practical. There is a study desk by the window, a single bed, and a shelf where I keep my books and class materials.

\textbf{Sample 2: Working Professional Persona (Band 7-8)}

My bedroom is simple and tidy, with a queen-sized bed and a small work corner. I try not to clutter it because I want it to be a place where I can properly relax.

\separator

\begin{ieltsquestion}
\textbf{What part of your home do you like the most?}
\end{ieltsquestion}
\textbf{Sample 1: Student Persona (Band 7-8)}

I like the desk area in my room the most. It is where I study, watch lectures, and sometimes have video calls with my family, so it feels like my own little corner.

\textbf{Sample 2: Working Professional Persona (Band 7-8)}

My favourite part is the balcony. It is not very big, but I can sit there with a cup of coffee in the morning and get some fresh air before work.

\separator

\begin{ieltsquestion}
\textbf{What's your favorite room in your apartment or house?}
\end{ieltsquestion}
\textbf{Sample 1: Student Persona (Band 7-8)}

My favourite room is probably my bedroom because it gives me privacy. Since I share the apartment with others, having a quiet personal space is really important.

\textbf{Sample 2: Working Professional Persona (Band 7-8)}

I would say the living room. It is where I unwind after work, watch a film, or call my friends, so it is the most relaxing part of the apartment.

\separator

\begin{ieltsquestion}
\textbf{What room does your family spend most of the time in?}
\end{ieltsquestion}
\textbf{Sample 1: Student Persona (Band 7-8)}

Back at my family home, we spend most of our time in the living room. My parents usually watch the news there, and I sit with them whenever I am home during the holidays.

\textbf{Sample 2: Working Professional Persona (Band 7-8)}

When my family comes over, we usually gather in the living room. It is the most comfortable space for chatting, eating snacks, and catching up with each other.

\separator

\begin{ieltsquestion}
\textbf{What do you usually do in your apartment?}
\end{ieltsquestion}
\textbf{Sample 1: Student Persona (Band 7-8)}

I usually study, cook simple meals, and rest there. On weekends, my roommates and I sometimes watch movies together or order food when we are too lazy to go out.

\textbf{Sample 2: Working Professional Persona (Band 7-8)}

I mostly relax, prepare meals, and do a bit of housework. Occasionally, I also work from home, so I need the apartment to be quiet enough for online meetings.

\separator

\begin{ieltsquestion}
\textbf{What makes you feel pleasant in your home?}
\end{ieltsquestion}
\textbf{Sample 1: Student Persona (Band 7-8)}

Small things make it pleasant, like clean bedding, enough sunlight, and a tidy desk. When the apartment is organised, I feel much less stressed about my studies.

\textbf{Sample 2: Working Professional Persona (Band 7-8)}

For me, it is the sense of calm. After dealing with traffic and deadlines, coming back to a clean, quiet space helps me mentally switch off.

\separator

\begin{ieltsquestion}
\textbf{How long have you lived there?}
\end{ieltsquestion}
\textbf{Sample 1: Student Persona (Band 7-8)}

I have lived there for almost a year. I moved in at the beginning of my second year because it was much closer to campus than my previous place.

\textbf{Sample 2: Working Professional Persona (Band 7-8)}

I have lived in my current apartment for about two years. I chose it mainly because it is close to work and the rent was reasonable when I first moved in.

\separator

\begin{ieltsquestion}
\textbf{Do you plan to live there for a long time?}
\end{ieltsquestion}
\textbf{Sample 1: Student Persona (Band 7-8)}

Probably not for a very long time. It is suitable while I am studying, but after graduation I may move to another area depending on where I find a job.

\textbf{Sample 2: Working Professional Persona (Band 7-8)}

I may stay there for another year or two, but not permanently. If my income increases, I would like to move to a slightly bigger place with a proper home office.

\separator

\begin{ieltsquestion}
\textbf{Are the transport facilities to your home very good?}
\end{ieltsquestion}
\textbf{Sample 1: Student Persona (Band 7-8)}

Yes, they are fairly good. There are bus stops nearby, and I can also ride my motorbike to university in about ten minutes, which is very convenient.

\textbf{Sample 2: Working Professional Persona (Band 7-8)}

They are decent, but traffic can still be frustrating during rush hour. The good thing is that my apartment is close enough to the office, so my commute is usually manageable.

\separator

\begin{ieltsquestion}
\textbf{What's the difference between where you are living now and where you have lived in the past?}
\end{ieltsquestion}
\textbf{Sample 1: Student Persona (Band 7-8)}

In the past, I lived with my family in a bigger house, so there was more space and less responsibility. Now, I live in a smaller shared apartment, but it has taught me to be more independent.

\textbf{Sample 2: Working Professional Persona (Band 7-8)}

I used to live with my parents, where everything felt more spacious and familiar. Now my apartment is smaller, but it is closer to work and gives me more freedom over my daily routine.

\separator

\begin{ieltsquestion}
\textbf{What kind of house or apartment do you want to live in in the future?}
\end{ieltsquestion}
\textbf{Sample 1: Student Persona (Band 7-8)}

In the future, I would like to live in a bright apartment with a balcony and a separate study area. It does not need to be huge, but it should feel peaceful and practical.

\textbf{Sample 2: Working Professional Persona (Band 7-8)}

I would love to own a modern apartment with two bedrooms, good soundproofing, and enough space for a home office. Since remote work is becoming more common, that would be extremely useful.

\separator

\begin{ieltsquestion}
\textbf{Do you think it is important to live in a comfortable environment?}
\end{ieltsquestion}
\textbf{Sample 1: Student Persona (Band 7-8)}

Yes, definitely. A comfortable living environment affects my mood and concentration, so it can make a real difference to how well I study.

\textbf{Sample 2: Working Professional Persona (Band 7-8)}

Absolutely. If your home is noisy, cramped, or messy, it is hard to recharge properly, and that can eventually affect your performance at work.

\separator


\vspace{0.8em}
\begin{center}
\renewcommand{\arraystretch}{1.65}
\rowcolors{2}{ForumFreshLight}{white}
\begin{tabularx}{\textwidth}{|V{0.22\textwidth}|C{0.07\textwidth}|I{0.20\textwidth}|Y|V{0.18\textwidth}|}
\hline
\rowcolor{ForumFreshBlue!20}
\textbf{Từ} & \textbf{Loại} & \textbf{IPA} & \textbf{English Meaning} & \textbf{Dịch nghĩa} \\
\hline
\textbf{rented apartment} & n phr & /ˈren.tɪd əˈpɑːt.mənt/ & a flat paid for monthly by a tenant & căn hộ thuê \\ \hline
\textbf{compact} & adj & /kəmˈpækt/ & small but arranged efficiently & nhỏ gọn \\ \hline
\textbf{low-maintenance} & adj & /ˌləʊ ˈmeɪn.tən.əns/ & requiring little time or effort to look after & ít tốn công chăm sóc \\ \hline
\textbf{natural light} & n phr & /ˈnætʃ.ər.əl laɪt/ & sunlight that enters a room & ánh sáng tự nhiên \\ \hline
\textbf{clutter} & v & /ˈklʌt.ər/ & to fill a space with too many things & làm bừa bộn \\ \hline
\textbf{unwind} & v & /ʌnˈwaɪnd/ & to relax after stress or work & thư giãn \\ \hline
\textbf{switch off} & phr v & /swɪtʃ ɒf/ & to stop thinking about work or stress & ngừng nghĩ về công việc/căng thẳng \\ \hline
\textbf{commute} & n & /kəˈmjuːt/ & the journey between home and work or school & quãng đường đi học/đi làm \\ \hline
\textbf{soundproofing} & n & /ˈsaʊndˌpruː.fɪŋ/ & material or design that blocks noise & khả năng cách âm \\ \hline
\textbf{recharge} & v & /ˌriːˈtʃɑːdʒ/ & to regain energy after being tired & nạp lại năng lượng \\ \hline
\end{tabularx}
\end{center}

\newpage

\section{The City You Live In}
\begin{ieltsquestion}
\textbf{What city do you live in?}
\end{ieltsquestion}
\textbf{Sample 1: Student Persona (Band 7-8)}

I currently live in Da Nang because my university is there. It is a coastal city, so life feels a bit more relaxed than in the really crowded urban centres.

\textbf{Sample 2: Working Professional Persona (Band 7-8)}

I live in Ho Chi Minh City at the moment. It is the biggest business hub in Vietnam, and I moved here mainly because there are more career opportunities in my field.

\separator

\begin{ieltsquestion}
\textbf{Do you like this city? Why?}
\end{ieltsquestion}
\textbf{Sample 1: Student Persona (Band 7-8)}

Yes, I like it a lot. It is lively enough for students, with cafes, beaches, and shopping areas, but it is still not too chaotic, so I can focus on my studies.

\textbf{Sample 2: Working Professional Persona (Band 7-8)}

I do, although it can be exhausting. The city is energetic and full of opportunities, so if you are ambitious and willing to work hard, it can be a very rewarding place to live.

\separator

\begin{ieltsquestion}
\textbf{Is this city your permanent residence?}
\end{ieltsquestion}
\textbf{Sample 1: Student Persona (Band 7-8)}

Not necessarily. I live here for university now, but I may move after graduation depending on where I can find a suitable job.

\textbf{Sample 2: Working Professional Persona (Band 7-8)}

For now, yes, but I am not sure whether it will be permanent. My career is here at the moment, but in the long run I might move somewhere quieter.

\separator

\begin{ieltsquestion}
\textbf{How long have you lived in this city?}
\end{ieltsquestion}
\textbf{Sample 1: Student Persona (Band 7-8)}

I have lived here for nearly two years. I moved here when I started university, and it took me a few months to get used to living away from home.

\textbf{Sample 2: Working Professional Persona (Band 7-8)}

I have lived in Ho Chi Minh City for about three years. At first, the traffic and pace of life were a bit overwhelming, but now I have adapted to them.

\separator

\begin{ieltsquestion}
\textbf{Are there big changes in this city?}
\end{ieltsquestion}
\textbf{Sample 1: Student Persona (Band 7-8)}

Yes, there are. New apartment buildings, cafes, and tourist services are appearing quite quickly, so the city feels more modern every year.

\textbf{Sample 2: Working Professional Persona (Band 7-8)}

Definitely. There are new office towers, shopping malls, and metro lines under development, although the rapid growth also brings problems like congestion and higher living costs.

\separator

\begin{ieltsquestion}
\textbf{What's the weather like where you live?}
\end{ieltsquestion}
\textbf{Sample 1: Student Persona (Band 7-8)}

It is usually warm and sunny, which is nice for going to the beach. However, the rainy season can be quite intense, and storms sometimes make travelling around difficult.

\textbf{Sample 2: Working Professional Persona (Band 7-8)}

Ho Chi Minh City is hot most of the year, with a clear rainy season. I do not mind the heat too much, but sudden downpours during rush hour can be really inconvenient.

\separator

\begin{ieltsquestion}
\textbf{Are the people friendly in the city?}
\end{ieltsquestion}
\textbf{Sample 1: Student Persona (Band 7-8)}

Yes, generally speaking. People in Da Nang are quite approachable, and I have found that shop owners, neighbours, and even strangers are usually willing to help.

\textbf{Sample 2: Working Professional Persona (Band 7-8)}

I would say people are friendly but also very busy. They may not stop for long conversations, but in everyday situations they are usually open, direct, and helpful.

\separator

\begin{ieltsquestion}
\textbf{Is the city friendly to children and old people?}
\end{ieltsquestion}
\textbf{Sample 1: Student Persona (Band 7-8)}

To some extent, yes. There are parks, beaches, and safe public spaces, but some roads are still quite busy, so parents and older people need to be careful when crossing streets.

\textbf{Sample 2: Working Professional Persona (Band 7-8)}

Not completely. There are good hospitals and services, but the traffic, noise, and crowded pavements can be difficult for children and elderly people.

\separator

\begin{ieltsquestion}
\textbf{Are there people of different ages living in this city?}
\end{ieltsquestion}
\textbf{Sample 1: Student Persona (Band 7-8)}

Yes, definitely. There are students, young workers, families, and retired people, especially because Da Nang is both a university city and a comfortable place to settle down.

\textbf{Sample 2: Working Professional Persona (Band 7-8)}

Yes, Ho Chi Minh City has people of all ages. You can see young professionals rushing to work, students hanging out after class, and older residents running small shops in local neighbourhoods.

\separator

\begin{ieltsquestion}
\textbf{Do you often see your neighbors?}
\end{ieltsquestion}
\textbf{Sample 1: Student Persona (Band 7-8)}

Yes, but mostly in passing. I see them in the hallway or near the entrance of the apartment building, and we usually just smile or say a quick hello.

\textbf{Sample 2: Working Professional Persona (Band 7-8)}

Not very often. I live in an apartment building where people keep to themselves, so I might see my neighbours in the lift, but we rarely have proper conversations.

\separator

\begin{ieltsquestion}
\textbf{Would you recommend your city to others?}
\end{ieltsquestion}
\textbf{Sample 1: Student Persona (Band 7-8)}

Yes, especially to students or people who want a balanced lifestyle. Da Nang is affordable, scenic, and not too stressful, which makes it a pleasant place to live.

\textbf{Sample 2: Working Professional Persona (Band 7-8)}

Yes, I would recommend Ho Chi Minh City to people who are career-oriented and enjoy a fast-paced lifestyle. It is not the easiest city to live in, but it offers a lot of room to grow.

\separator


\vspace{0.8em}
\begin{center}
\renewcommand{\arraystretch}{1.65}
\rowcolors{2}{ForumFreshLight}{white}
\begin{tabularx}{\textwidth}{|V{0.22\textwidth}|C{0.07\textwidth}|I{0.20\textwidth}|Y|V{0.18\textwidth}|}
\hline
\rowcolor{ForumFreshBlue!20}
\textbf{Từ} & \textbf{Loại} & \textbf{IPA} & \textbf{English Meaning} & \textbf{Dịch nghĩa} \\
\hline
\textbf{urban centre} & n phr & /ˈɜː.bən ˈsen.tər/ & a busy central city area & trung tâm đô thị \\ \hline
\textbf{business hub} & n phr & /ˈbɪz.nəs hʌb/ & an important centre for business activity & trung tâm kinh doanh \\ \hline
\textbf{chaotic} & adj & /keɪˈɒt.ɪk/ & confused, busy, and lacking order & hỗn loạn \\ \hline
\textbf{rewarding} & adj & /rɪˈwɔː.dɪŋ/ & giving satisfaction or benefits & đáng giá, thỏa mãn \\ \hline
\textbf{overwhelming} & adj & /ˌəʊ.vəˈwel.mɪŋ/ & too intense or difficult to deal with & choáng ngợp \\ \hline
\textbf{congestion} & n & /kənˈdʒes.tʃən/ & overcrowding, especially traffic & sự tắc nghẽn \\ \hline
\textbf{downpour} & n & /ˈdaʊn.pɔːr/ & a heavy fall of rain & trận mưa lớn \\ \hline
\textbf{approachable} & adj & /əˈprəʊ.tʃə.bəl/ & friendly and easy to talk to & dễ gần \\ \hline
\textbf{keep to oneself} & idiom & /kiːp tuː wʌnˈself/ & to avoid much social interaction & sống khép kín, ít giao tiếp \\ \hline
\textbf{career-oriented} & adj & /kəˈrɪə ˌɔː.ri.en.tɪd/ & focused on career development & định hướng sự nghiệp \\ \hline
\end{tabularx}
\end{center}

\newpage

\section{The Area You Live In}
\begin{ieltsquestion}
\textbf{Do you live in a noisy or a quiet area?}
\end{ieltsquestion}
\textbf{Sample 1: Student Persona (Band 7-8)}

I would say it is fairly quiet, especially at night. During the day, there is some traffic because many students live nearby, but it is not noisy enough to disturb my study routine.

\textbf{Sample 2: Working Professional Persona (Band 7-8)}

My area is a bit noisy because it is close to a main road and several office buildings. Still, my apartment is on a higher floor, so the noise is manageable most of the time.

\separator

\begin{ieltsquestion}
\textbf{Do you like the area that you live in?}
\end{ieltsquestion}
\textbf{Sample 1: Student Persona (Band 7-8)}

Yes, I do. It is close to my university, affordable food stalls, and a few cafes, so it is very convenient for student life.

\textbf{Sample 2: Working Professional Persona (Band 7-8)}

Yes, overall I like it. It is not the most peaceful neighbourhood, but it is practical because I can get to work quickly and run errands without travelling too far.

\separator

\begin{ieltsquestion}
\textbf{Where do you like to go in that area?}
\end{ieltsquestion}
\textbf{Sample 1: Student Persona (Band 7-8)}

I often go to a small coffee shop near my apartment. It has a quiet upstairs space, so I usually study there when I need a change of scenery.

\textbf{Sample 2: Working Professional Persona (Band 7-8)}

I usually go to a nearby convenience store or a small gym after work. On weekends, I sometimes visit a brunch place in the area because it helps me unwind.

\separator

\begin{ieltsquestion}
\textbf{Are the people in your neighborhood nice and friendly?}
\end{ieltsquestion}
\textbf{Sample 1: Student Persona (Band 7-8)}

Yes, most of them are quite friendly. The food stall owners near my building even remember what I usually order, which makes the area feel more welcoming.

\textbf{Sample 2: Working Professional Persona (Band 7-8)}

They are polite, though people tend to keep to themselves. In a busy city, I think that is normal, but whenever I need help, the security guards and shopkeepers are approachable.

\separator

\begin{ieltsquestion}
\textbf{Do you know any of your neighbours?}
\end{ieltsquestion}
\textbf{Sample 1: Student Persona (Band 7-8)}

I know a few students who live on the same floor. We are not extremely close, but we sometimes chat in the hallway or lend each other small things like chargers.

\textbf{Sample 2: Working Professional Persona (Band 7-8)}

Only a few, to be honest. I recognise the people who live next door and we greet each other in the lift, but our conversations are usually very brief.

\separator

\begin{ieltsquestion}
\textbf{Do you know any famous people in your area?}
\end{ieltsquestion}
\textbf{Sample 1: Student Persona (Band 7-8)}

Not personally. I have heard that a few local influencers sometimes film videos at cafes nearby, but I have never actually met anyone famous there.

\textbf{Sample 2: Working Professional Persona (Band 7-8)}

No, I do not know any famous people in my area. There are some successful business owners around, but they are not celebrities in the usual sense.

\separator

\begin{ieltsquestion}
\textbf{What are some changes in the area recently?}
\end{ieltsquestion}
\textbf{Sample 1: Student Persona (Band 7-8)}

Recently, several new cafes and convenience stores have opened because more students are moving into the area. It is becoming livelier, but luckily it still feels quite safe.

\textbf{Sample 2: Working Professional Persona (Band 7-8)}

There has been a lot of construction recently, especially new apartment blocks and office spaces. It makes the area more modern, but it also causes dust and traffic jams from time to time.

\separator


\vspace{0.8em}
\begin{center}
\renewcommand{\arraystretch}{1.65}
\rowcolors{2}{ForumFreshLight}{white}
\begin{tabularx}{\textwidth}{|V{0.22\textwidth}|C{0.07\textwidth}|I{0.20\textwidth}|Y|V{0.18\textwidth}|}
\hline
\rowcolor{ForumFreshBlue!20}
\textbf{Từ} & \textbf{Loại} & \textbf{IPA} & \textbf{English Meaning} & \textbf{Dịch nghĩa} \\
\hline
\textbf{disturb} & v & /dɪˈstɜːb/ & to interrupt or make someone lose focus & làm phiền, gây gián đoạn \\ \hline
\textbf{manageable} & adj & /ˈmæn.ɪ.dʒə.bəl/ & possible to deal with without too much difficulty & có thể kiểm soát được \\ \hline
\textbf{run errands} & v phr & /rʌn ˈer.əndz/ & to do small necessary tasks outside the home & đi làm việc vặt \\ \hline
\textbf{change of scenery} & n phr & /tʃeɪndʒ əv ˈsiː.nər.i/ & a different place or environment & thay đổi không gian \\ \hline
\textbf{unwind} & v & /ʌnˈwaɪnd/ & to relax after stress or work & thư giãn \\ \hline
\textbf{welcoming} & adj & /ˈwel.kəm.ɪŋ/ & friendly and pleasant & thân thiện, dễ chịu \\ \hline
\textbf{keep to themselves} & idiom & /kiːp tuː ðəmˈselvz/ & to avoid much social interaction & sống khép kín \\ \hline
\textbf{approachable} & adj & /əˈprəʊ.tʃə.bəl/ & easy and pleasant to talk to & dễ gần \\ \hline
\textbf{lively} & adj & /ˈlaɪv.li/ & full of activity and energy & sôi động \\ \hline
\textbf{apartment block} & n phr & /əˈpɑːt.mənt blɒk/ & a large building containing many apartments & tòa chung cư \\ \hline
\end{tabularx}
\end{center}

\newpage

\section*{PART 1: DAILY LIFE \& HABITS}
\addcontentsline{toc}{section}{PART 1: DAILY LIFE \& HABITS}

\section{Morning Time}
\begin{ieltsquestion}
\textbf{Do you like getting up early in the morning?}
\end{ieltsquestion}
\textbf{Sample 1: Student Persona (Band 7-8)}

To be honest, I am not naturally a morning person. I like the feeling of waking up early, but during exam periods I often stay up late, so getting out of bed can be a real struggle.

\textbf{Sample 2: Working Professional Persona (Band 7-8)}

Yes, I do, mainly because my workday starts quite early. If I wake up ahead of time, I can have breakfast properly and avoid rushing, which puts me in a better mood.

\separator

\begin{ieltsquestion}
\textbf{What do you usually do in the morning?}
\end{ieltsquestion}
\textbf{Sample 1: Student Persona (Band 7-8)}

On weekdays, I usually wash up, grab a quick breakfast, and check my class schedule. If I have an early lecture, I leave home in a hurry, but when my class starts later, I review my notes for a while.

\textbf{Sample 2: Working Professional Persona (Band 7-8)}

I normally make coffee, read a few work messages, and plan my priorities for the day. After that, I commute to the office, and I usually listen to a podcast on the way.

\separator

\begin{ieltsquestion}
\textbf{What did you do in the morning when you were little? Why?}
\end{ieltsquestion}
\textbf{Sample 1: Student Persona (Band 7-8)}

When I was little, my mornings were much more structured. My parents woke me up, prepared breakfast, and made sure I packed my school bag because I was not responsible enough to organise everything myself.

\textbf{Sample 2: Working Professional Persona (Band 7-8)}

As a child, I usually watched cartoons while eating breakfast before school. It was a small routine that made mornings less boring, especially because I did not really enjoy getting up early back then.

\separator

\begin{ieltsquestion}
\textbf{Are there any differences between what you do in the morning now and what you did in the past?}
\end{ieltsquestion}
\textbf{Sample 1: Student Persona (Band 7-8)}

Yes, definitely. In the past, my parents handled most things for me, but now I have to manage my own time, cook simple meals, and make sure I am not late for class.

\textbf{Sample 2: Working Professional Persona (Band 7-8)}

There is a big difference. My mornings used to be quite carefree, but now they are more purposeful because I need to prepare mentally for meetings, deadlines, and decisions at work.

\separator

\begin{ieltsquestion}
\textbf{Do you spend your mornings doing the same things on both weekends and weekdays? Why?}
\end{ieltsquestion}
\textbf{Sample 1: Student Persona (Band 7-8)}

No, not really. On weekdays, my mornings are shaped by my timetable, but on weekends I usually sleep in, have a slower breakfast, and maybe do some light revision if I have exams coming up.

\textbf{Sample 2: Working Professional Persona (Band 7-8)}

No, my weekend mornings are much more relaxed. I do not check work emails unless something urgent comes up, and I often go for a walk or have brunch with friends.

\separator


\vspace{0.8em}
\begin{center}
\renewcommand{\arraystretch}{1.65}
\rowcolors{2}{ForumFreshLight}{white}
\begin{tabularx}{\textwidth}{|V{0.22\textwidth}|C{0.07\textwidth}|I{0.20\textwidth}|Y|V{0.18\textwidth}|}
\hline
\rowcolor{ForumFreshBlue!20}
\textbf{Từ} & \textbf{Loại} & \textbf{IPA} & \textbf{English Meaning} & \textbf{Dịch nghĩa} \\
\hline
\textbf{morning person} & n phr & /ˈmɔː.nɪŋ ˈpɜː.sən/ & someone who feels energetic in the morning & người thích buổi sáng/dậy sớm \\ \hline
\textbf{exam period} & n phr & /ɪɡˈzæm ˈpɪə.ri.əd/ & the time when students take exams & giai đoạn thi cử \\ \hline
\textbf{ahead of time} & adv phr & /əˈhed əv taɪm/ & earlier than necessary or expected & sớm hơn dự kiến \\ \hline
\textbf{put someone in a better mood} & v phr & /pʊt ˈsʌm.wʌn ɪn ə ˈbet.ə muːd/ & to make someone feel happier & khiến ai có tâm trạng tốt hơn \\ \hline
\textbf{structured} & adj & /ˈstrʌk.tʃəd/ & organised and planned clearly & có cấu trúc, có kế hoạch \\ \hline
\textbf{routine} & n & /ruːˈtiːn/ & a regular way of doing things & thói quen, lịch trình \\ \hline
\textbf{carefree} & adj & /ˈkeə.friː/ & free from worries or responsibilities & vô tư, không lo nghĩ \\ \hline
\textbf{purposeful} & adj & /ˈpɜː.pəs.fəl/ & having a clear aim or intention & có mục đích rõ ràng \\ \hline
\textbf{sleep in} & phr v & /sliːp ɪn/ & to sleep later than usual & ngủ nướng \\ \hline
\textbf{brunch} & n & /brʌntʃ/ & a meal between breakfast and lunch & bữa sáng muộn \\ \hline
\end{tabularx}
\end{center}

\newpage

\section{Spare Time}
\begin{ieltsquestion}
\textbf{Do you often have free time?}
\end{ieltsquestion}
\textbf{Sample 1: Student Persona (Band 7-8)}

I have some free time, but not as much as I would like. My schedule changes from week to week, and assignments can easily take over my evenings.

\textbf{Sample 2: Working Professional Persona (Band 7-8)}

Not really during weekdays. After work, commuting, and housework, I usually only have a small window of time to relax before going to bed.

\separator

\begin{ieltsquestion}
\textbf{What do you usually do in your spare time?}
\end{ieltsquestion}
\textbf{Sample 1: Student Persona (Band 7-8)}

I usually watch videos, listen to music, or meet friends at a cafe near campus. If I have more time, I like doing a bit of graphic design because it feels creative and relaxing.

\textbf{Sample 2: Working Professional Persona (Band 7-8)}

I usually unwind by going to the gym or watching a light series. Sometimes I cook something simple because it helps me slow down after a demanding day.

\separator

\begin{ieltsquestion}
\textbf{Do you have more free time on weekdays or at weekends?}
\end{ieltsquestion}
\textbf{Sample 1: Student Persona (Band 7-8)}

Definitely at weekends. On weekdays, classes and homework take up most of my time, but weekends give me more breathing space to rest or catch up with friends.

\textbf{Sample 2: Working Professional Persona (Band 7-8)}

I have much more free time at weekends. Weekdays are fairly packed, so Saturday and Sunday are when I can properly recharge and do things at my own pace.

\separator

\begin{ieltsquestion}
\textbf{Which day do you have more free time on, Saturday or Sunday?}
\end{ieltsquestion}
\textbf{Sample 1: Student Persona (Band 7-8)}

Usually Sunday. I sometimes have group meetings or part-time tasks on Saturday, while Sunday is the day when I can stay home, tidy my room, and prepare for the next week.

\textbf{Sample 2: Working Professional Persona (Band 7-8)}

I would say Sunday. Saturday is often filled with errands or social plans, but Sunday feels quieter, so I can sleep a bit longer and enjoy a slower day.

\separator

\begin{ieltsquestion}
\textbf{Would you like to have more free time in the future?}
\end{ieltsquestion}
\textbf{Sample 1: Student Persona (Band 7-8)}

Yes, absolutely. I would like to have more time for hobbies and exercise because at the moment I sometimes feel that my life revolves too much around deadlines.

\textbf{Sample 2: Working Professional Persona (Band 7-8)}

Yes, I would. I do not want endless free time, but I would like a healthier balance where I can pursue personal interests without feeling guilty about unfinished work.

\separator


\vspace{0.8em}
\begin{center}
\renewcommand{\arraystretch}{1.65}
\rowcolors{2}{ForumFreshLight}{white}
\begin{tabularx}{\textwidth}{|V{0.22\textwidth}|C{0.07\textwidth}|I{0.20\textwidth}|Y|V{0.18\textwidth}|}
\hline
\rowcolor{ForumFreshBlue!20}
\textbf{Từ} & \textbf{Loại} & \textbf{IPA} & \textbf{English Meaning} & \textbf{Dịch nghĩa} \\
\hline
\textbf{take over} & phr v & /teɪk ˈəʊ.vər/ & to occupy or control something completely & chiếm hết, lấn át \\ \hline
\textbf{small window of time} & n phr & /smɔːl ˈwɪn.dəʊ əv taɪm/ & a short available period & khoảng thời gian ngắn \\ \hline
\textbf{unwind} & v & /ʌnˈwaɪnd/ & to relax after stress or work & thư giãn \\ \hline
\textbf{demanding} & adj & /dɪˈmɑːn.dɪŋ/ & requiring a lot of effort & đòi hỏi nhiều sức lực \\ \hline
\textbf{breathing space} & n phr & /ˈbriː.ðɪŋ speɪs/ & time or room to rest and think & khoảng nghỉ, không gian để thở \\ \hline
\textbf{recharge} & v & /ˌriːˈtʃɑːdʒ/ & to regain energy & nạp lại năng lượng \\ \hline
\textbf{at my own pace} & adv phr & /æt maɪ əʊn peɪs/ & at a speed that suits me & theo nhịp độ của riêng mình \\ \hline
\textbf{run errands} & v phr & /rʌn ˈer.əndz/ & to do small necessary tasks & làm việc vặt \\ \hline
\textbf{revolve around} & phr v & /rɪˈvɒlv əˈraʊnd/ & to have something as the main focus & xoay quanh \\ \hline
\textbf{unfinished work} & n phr & /ʌnˈfɪn.ɪʃt wɜːk/ & tasks that have not been completed & công việc chưa hoàn thành \\ \hline
\end{tabularx}
\end{center}

\newpage

\section{Day Off}
\begin{ieltsquestion}
\textbf{When was the last time you had a few days off?}
\end{ieltsquestion}
\textbf{Sample 1: Student Persona (Band 7-8)}

The last time was after my final exams last semester. I had almost a week without classes, so I went back home, slept a lot, and spent some proper time with my family.

\textbf{Sample 2: Working Professional Persona (Band 7-8)}

I had a few days off during the last public holiday. I did not travel anywhere far, but I used the time to rest, sort out personal errands, and catch up on sleep.

\separator

\begin{ieltsquestion}
\textbf{What do you usually do when you have days off?}
\end{ieltsquestion}
\textbf{Sample 1: Student Persona (Band 7-8)}

I usually slow down and do things I cannot fit into a normal school week. I might meet friends, watch a series, or simply stay at home and recharge.

\textbf{Sample 2: Working Professional Persona (Band 7-8)}

On my days off, I try not to think too much about work. I usually clean my apartment, go to a cafe, exercise a little, and maybe have dinner with friends in the evening.

\separator

\begin{ieltsquestion}
\textbf{What would you like to do if you had a day off tomorrow?}
\end{ieltsquestion}
\textbf{Sample 1: Student Persona (Band 7-8)}

If I had a day off tomorrow, I would probably go to the beach in the morning and then study lightly in the afternoon. That would give me a nice balance between relaxation and responsibility.

\textbf{Sample 2: Working Professional Persona (Band 7-8)}

I would love to have a slow morning, make breakfast at home, and then get a massage or go for a swim. I think one quiet day like that would help me reset mentally.

\separator

\begin{ieltsquestion}
\textbf{Do you usually spend your days off with your parents or with your friends?}
\end{ieltsquestion}
\textbf{Sample 1: Student Persona (Band 7-8)}

It depends. If I am back in my hometown, I spend most of my time with my parents, but during the semester I usually hang out with friends because we live close to each other.

\textbf{Sample 2: Working Professional Persona (Band 7-8)}

I spend them with both, but probably more often with friends because my parents live in another city. Still, I call my family regularly, especially when I have a more relaxed day.

\separator


\vspace{0.8em}
\begin{center}
\renewcommand{\arraystretch}{1.65}
\rowcolors{2}{ForumFreshLight}{white}
\begin{tabularx}{\textwidth}{|V{0.22\textwidth}|C{0.07\textwidth}|I{0.20\textwidth}|Y|V{0.18\textwidth}|}
\hline
\rowcolor{ForumFreshBlue!20}
\textbf{Từ} & \textbf{Loại} & \textbf{IPA} & \textbf{English Meaning} & \textbf{Dịch nghĩa} \\
\hline
\textbf{final exams} & n phr & /ˈfaɪ.nəl ɪɡˈzæmz/ & important exams at the end of a term & kỳ thi cuối kỳ \\ \hline
\textbf{public holiday} & n phr & /ˌpʌb.lɪk ˈhɒl.ə.deɪ/ & a national day off work or school & ngày nghỉ lễ \\ \hline
\textbf{sort out} & phr v & /sɔːt aʊt/ & to organise or deal with something & sắp xếp, xử lý \\ \hline
\textbf{catch up on sleep} & v phr & /kætʃ ʌp ɒn sliːp/ & to sleep more after not sleeping enough & ngủ bù \\ \hline
\textbf{slow down} & phr v & /sləʊ daʊn/ & to become less busy or active & chậm lại, nghỉ nhịp \\ \hline
\textbf{recharge} & v & /ˌriːˈtʃɑːdʒ/ & to regain energy & nạp lại năng lượng \\ \hline
\textbf{reset mentally} & v phr & /ˌriːˈset ˈmen.təl.i/ & to refresh one's mind & làm mới tinh thần \\ \hline
\textbf{hang out} & phr v & /hæŋ aʊt/ & to spend relaxed time with others & đi chơi, tụ tập \\ \hline
\end{tabularx}
\end{center}

\newpage

\section{Having a Break}
\begin{ieltsquestion}
\textbf{How often do you take a rest or a break?}
\end{ieltsquestion}
\textbf{Sample 1: Student Persona (Band 7-8)}

I try to take a short break every hour when I am studying. If I study for too long without stopping, my concentration drops and I end up reading the same paragraph again and again.

\textbf{Sample 2: Working Professional Persona (Band 7-8)}

I usually take two or three short breaks during the workday. They are not long, maybe just five or ten minutes, but they help me stay focused and avoid feeling drained.

\separator

\begin{ieltsquestion}
\textbf{What do you usually do when you are resting?}
\end{ieltsquestion}
\textbf{Sample 1: Student Persona (Band 7-8)}

I usually stretch, drink water, or scroll through my phone for a few minutes. Sometimes I step outside for fresh air, especially when I have been staring at my laptop for too long.

\textbf{Sample 2: Working Professional Persona (Band 7-8)}

I often walk around the office, make a cup of tea, or chat briefly with a colleague. I try not to open social media too much because it can easily turn a short break into a long distraction.

\separator

\begin{ieltsquestion}
\textbf{Do you take a nap when you are taking your rest?}
\end{ieltsquestion}
\textbf{Sample 1: Student Persona (Band 7-8)}

Yes, sometimes, especially after lunch or during exam weeks. A short nap of around twenty minutes can make me feel much more alert without making me too sleepy afterwards.

\textbf{Sample 2: Working Professional Persona (Band 7-8)}

Not very often on workdays because there is no suitable place for it. But on weekends, I do enjoy taking a quick nap in the afternoon if I feel physically tired.

\separator

\begin{ieltsquestion}
\textbf{How do you feel after taking a nap?}
\end{ieltsquestion}
\textbf{Sample 1: Student Persona (Band 7-8)}

If the nap is short, I feel refreshed and ready to study again. But if I sleep for too long, I wake up groggy and it actually becomes harder to concentrate.

\textbf{Sample 2: Working Professional Persona (Band 7-8)}

A good nap makes me feel recharged and calmer. It is almost like pressing a reset button, especially after a busy morning or a stressful week.

\separator


\vspace{0.8em}
\begin{center}
\renewcommand{\arraystretch}{1.65}
\rowcolors{2}{ForumFreshLight}{white}
\begin{tabularx}{\textwidth}{|V{0.22\textwidth}|C{0.07\textwidth}|I{0.20\textwidth}|Y|V{0.18\textwidth}|}
\hline
\rowcolor{ForumFreshBlue!20}
\textbf{Từ} & \textbf{Loại} & \textbf{IPA} & \textbf{English Meaning} & \textbf{Dịch nghĩa} \\
\hline
\textbf{concentration} & n & /ˌkɒn.sənˈtreɪ.ʃən/ & the ability to focus & sự tập trung \\ \hline
\textbf{end up} & phr v & /end ʌp/ & to finally be in a situation & rốt cuộc là \\ \hline
\textbf{drained} & adj & /dreɪnd/ & very tired and low on energy & kiệt sức \\ \hline
\textbf{stretch} & v & /stretʃ/ & to extend the body or muscles & duỗi người, giãn cơ \\ \hline
\textbf{distraction} & n & /dɪˈstræk.ʃən/ & something that takes attention away & sự xao nhãng \\ \hline
\textbf{alert} & adj & /əˈlɜːt/ & awake and able to think clearly & tỉnh táo \\ \hline
\textbf{physically tired} & adj phr & /ˈfɪz.ɪ.kəl.i ˈtaɪəd/ & tired in the body & mệt về thể chất \\ \hline
\textbf{refreshed} & adj & /rɪˈfreʃt/ & feeling rested and energetic again & sảng khoái \\ \hline
\textbf{groggy} & adj & /ˈɡrɒɡ.i/ & weak or confused after sleep & uể oải sau khi ngủ \\ \hline
\textbf{press a reset button} & idiom & /pres ə ˈriː.set ˈbʌt.ən/ & to start again with a clearer mind & như khởi động lại tinh thần \\ \hline
\end{tabularx}
\end{center}

\newpage

\section{Going Out}
\begin{ieltsquestion}
\textbf{Do you always take your mobile phone with you when going out?}
\end{ieltsquestion}
\textbf{Sample 1: Student Persona (Band 7-8)}

Yes, almost always. I use my phone for maps, messages, online payments, and even checking my class group chat, so leaving it at home would make me feel a bit disconnected.

\textbf{Sample 2: Working Professional Persona (Band 7-8)}

Yes, definitely. My phone is essential because I need it for work calls, banking apps, ride-hailing services, and keeping in touch with colleagues or clients.

\separator

\begin{ieltsquestion}
\textbf{Do you often bring cash with you?}
\end{ieltsquestion}
\textbf{Sample 1: Student Persona (Band 7-8)}

I usually bring a small amount of cash, just in case. Most places accept bank transfers now, but some food stalls near my university still prefer cash.

\textbf{Sample 2: Working Professional Persona (Band 7-8)}

Not much, to be honest. I mainly pay by card or e-wallet, but I still keep some cash in my wallet for parking fees or small street-food places.

\separator

\begin{ieltsquestion}
\textbf{How often do you use cash?}
\end{ieltsquestion}
\textbf{Sample 1: Student Persona (Band 7-8)}

I use cash a few times a week. It is mostly for inexpensive things like breakfast, photocopying documents, or buying snacks from small vendors.

\textbf{Sample 2: Working Professional Persona (Band 7-8)}

I use cash less and less these days. Digital payments are faster and easier to track, so cash is more of a backup option for me now.

\separator

\begin{ieltsquestion}
\textbf{Do you bring food or snacks with you when going out?}
\end{ieltsquestion}
\textbf{Sample 1: Student Persona (Band 7-8)}

Sometimes I do, especially when I have a long day at university. I might bring a sandwich or some biscuits because buying food on campus every day can add up.

\textbf{Sample 2: Working Professional Persona (Band 7-8)}

Only occasionally. If I know I will be stuck in meetings or traffic, I may carry a protein bar or some nuts, but usually I just buy something nearby.

\separator


\vspace{0.8em}
\begin{center}
\renewcommand{\arraystretch}{1.65}
\rowcolors{2}{ForumFreshLight}{white}
\begin{tabularx}{\textwidth}{|V{0.22\textwidth}|C{0.07\textwidth}|I{0.20\textwidth}|Y|V{0.18\textwidth}|}
\hline
\rowcolor{ForumFreshBlue!20}
\textbf{Từ} & \textbf{Loại} & \textbf{IPA} & \textbf{English Meaning} & \textbf{Dịch nghĩa} \\
\hline
\textbf{disconnected} & adj & /ˌdɪs.kəˈnek.tɪd/ & not connected or out of touch & mất kết nối \\ \hline
\textbf{essential} & adj & /ɪˈsen.ʃəl/ & extremely important or necessary & thiết yếu \\ \hline
\textbf{ride-hailing service} & n phr & /ˈraɪdˌheɪ.lɪŋ ˈsɜː.vɪs/ & an app-based transport service & dịch vụ gọi xe công nghệ \\ \hline
\textbf{just in case} & adv phr & /dʒʌst ɪn keɪs/ & as a precaution & phòng khi cần \\ \hline
\textbf{e-wallet} & n & /ˈiː ˌwɒl.ɪt/ & a digital wallet for payments & ví điện tử \\ \hline
\textbf{vendor} & n & /ˈven.dər/ & a person who sells things & người bán hàng \\ \hline
\textbf{backup option} & n phr & /ˈbæk.ʌp ˈɒp.ʃən/ & an alternative used if needed & phương án dự phòng \\ \hline
\textbf{add up} & phr v & /æd ʌp/ & to gradually become a large amount & cộng dồn thành nhiều \\ \hline
\textbf{protein bar} & n phr & /ˈprəʊ.tiːn bɑːr/ & a snack bar high in protein & thanh đồ ăn giàu đạm \\ \hline
\end{tabularx}
\end{center}

\newpage

\section{Walking}
\begin{ieltsquestion}
\textbf{Do you walk a lot?}
\end{ieltsquestion}
\textbf{Sample 1: Student Persona (Band 7-8)}

Yes, quite a lot, especially around campus. My classrooms, library, and favourite cafe are all within walking distance, so walking is part of my daily routine.

\textbf{Sample 2: Working Professional Persona (Band 7-8)}

Not as much as I should. I sit at a desk for most of the day, so I try to walk after work or during lunch breaks to stay a bit more active.

\separator

\begin{ieltsquestion}
\textbf{Where did you go for a walk lately?}
\end{ieltsquestion}
\textbf{Sample 1: Student Persona (Band 7-8)}

A few days ago, I walked along the beach with one of my classmates after our evening lecture. It was breezy and quiet, so it helped me clear my head.

\textbf{Sample 2: Working Professional Persona (Band 7-8)}

Recently, I went for a walk around a park near my apartment. I only walked for about twenty minutes, but it was enough to loosen up after sitting in meetings all day.

\separator

\begin{ieltsquestion}
\textbf{Did you often go outside to have a walk when you were a child?}
\end{ieltsquestion}
\textbf{Sample 1: Student Persona (Band 7-8)}

Yes, I did. When I was little, my parents often took me for walks around the neighbourhood after dinner because it was a simple way for the family to spend time together.

\textbf{Sample 2: Working Professional Persona (Band 7-8)}

Yes, quite often. I grew up in a quieter area, so I used to walk or ride a bike with other kids in the neighbourhood almost every afternoon.

\separator

\begin{ieltsquestion}
\textbf{Why do people like to walk in parks?}
\end{ieltsquestion}
\textbf{Sample 1: Student Persona (Band 7-8)}

I think parks give people a break from screens and traffic. The fresh air, trees, and open space make walking feel more relaxing than just walking along a busy road.

\textbf{Sample 2: Working Professional Persona (Band 7-8)}

People like parks because they offer a rare sense of calm in crowded cities. Walking there can reduce stress, improve mood, and make people feel healthier without needing expensive equipment.

\separator

\begin{ieltsquestion}
\textbf{Where would you like to take a long walk if you had the chance?}
\end{ieltsquestion}
\textbf{Sample 1: Student Persona (Band 7-8)}

I would love to take a long walk along a coastal road at sunset. I think the sea breeze and the view would make it feel peaceful rather than tiring.

\textbf{Sample 2: Working Professional Persona (Band 7-8)}

If I had the chance, I would take a long walk in a quiet mountain town like Da Lat. The cool weather and slower pace would be a nice escape from city life.

\separator


\vspace{0.8em}
\begin{center}
\renewcommand{\arraystretch}{1.65}
\rowcolors{2}{ForumFreshLight}{white}
\begin{tabularx}{\textwidth}{|V{0.22\textwidth}|C{0.07\textwidth}|I{0.20\textwidth}|Y|V{0.18\textwidth}|}
\hline
\rowcolor{ForumFreshBlue!20}
\textbf{Từ} & \textbf{Loại} & \textbf{IPA} & \textbf{English Meaning} & \textbf{Dịch nghĩa} \\
\hline
\textbf{within walking distance} & adv phr & /wɪˈðɪn ˈwɔː.kɪŋ ˈdɪs.təns/ & close enough to walk to & trong khoảng cách đi bộ \\ \hline
\textbf{daily routine} & n phr & /ˈdeɪ.li ruːˈtiːn/ & regular activities done every day & thói quen hằng ngày \\ \hline
\textbf{breezy} & adj & /ˈbriː.zi/ & pleasantly windy & có gió nhẹ, thoáng mát \\ \hline
\textbf{clear my head} & v phr & /klɪər maɪ hed/ & to feel mentally calmer and clearer & làm đầu óc tỉnh táo \\ \hline
\textbf{loosen up} & phr v & /ˈluː.sən ʌp/ & to relax tight muscles or tension & thư giãn, giãn cơ \\ \hline
\textbf{open space} & n phr & /ˈəʊ.pən speɪs/ & an outdoor area without many buildings & không gian mở \\ \hline
\textbf{sense of calm} & n phr & /sens əv kɑːm/ & a feeling of peace & cảm giác bình yên \\ \hline
\textbf{coastal road} & n phr & /ˈkəʊ.stəl rəʊd/ & a road near or along the sea & đường ven biển \\ \hline
\textbf{sea breeze} & n phr & /siː briːz/ & a light wind from the sea & gió biển \\ \hline
\textbf{escape from city life} & n phr & /ɪˈskeɪp frəm ˈsɪt.i laɪf/ & a break from urban pressure & thoát khỏi nhịp sống thành phố \\ \hline
\end{tabularx}
\end{center}

\newpage

\section*{PART 1: PERSONAL PREFERENCES \& ACTIVITIES}
\addcontentsline{toc}{section}{PART 1: PERSONAL PREFERENCES \& ACTIVITIES}

\section{Hobby}
\begin{ieltsquestion}
\textbf{Do you have any hobbies?}
\end{ieltsquestion}
\textbf{Sample 1: Student Persona (Band 7-8)}

Yes, my main hobby is graphic design. I started doing it for fun, but now I sometimes design posters for student clubs, so it feels both enjoyable and useful.

\textbf{Sample 2: Working Professional Persona (Band 7-8)}

Yes, I like cooking and going to the gym. Cooking helps me relax after work, while exercising keeps me from feeling too stiff after sitting at a desk all day.

\separator

\begin{ieltsquestion}
\textbf{Did you have any hobbies when you were a child?}
\end{ieltsquestion}
\textbf{Sample 1: Student Persona (Band 7-8)}

Yes, I used to draw a lot when I was little. I could spend hours copying cartoon characters, and I think that was where my interest in design first came from.

\textbf{Sample 2: Working Professional Persona (Band 7-8)}

Definitely. When I was a child, I loved collecting stickers and playing badminton with neighbours. They were simple hobbies, but they made my childhood feel very active and social.

\separator

\begin{ieltsquestion}
\textbf{Do you have a hobby that you've had since childhood?}
\end{ieltsquestion}
\textbf{Sample 1: Student Persona (Band 7-8)}

Yes, drawing is the one hobby that has stayed with me. The tools have changed from coloured pencils to design software, but the basic interest is still the same.

\textbf{Sample 2: Working Professional Persona (Band 7-8)}

Yes, I still enjoy badminton, although I do not play as often as before. Whenever I have time, I book a court with friends because it reminds me of my childhood.

\separator

\begin{ieltsquestion}
\textbf{Do you have the same hobbies as your family members?}
\end{ieltsquestion}
\textbf{Sample 1: Student Persona (Band 7-8)}

Not exactly. My parents prefer gardening and watching TV, while I am more into creative and digital activities, but we all enjoy listening to music at home.

\textbf{Sample 2: Working Professional Persona (Band 7-8)}

Partly, yes. My mother enjoys cooking, and I think I picked up that habit from her, although I prefer trying quick modern recipes rather than traditional dishes that take hours.

\separator


\vspace{0.8em}
\begin{center}
\renewcommand{\arraystretch}{1.65}
\rowcolors{2}{ForumFreshLight}{white}
\begin{tabularx}{\textwidth}{|V{0.22\textwidth}|C{0.07\textwidth}|I{0.20\textwidth}|Y|V{0.18\textwidth}|}
\hline
\rowcolor{ForumFreshBlue!20}
\textbf{Từ} & \textbf{Loại} & \textbf{IPA} & \textbf{English Meaning} & \textbf{Dịch nghĩa} \\
\hline
\textbf{graphic design} & n phr & /ˈɡræf.ɪk dɪˈzaɪn/ & the art of creating visual materials & thiết kế đồ họa \\ \hline
\textbf{for fun} & adv phr & /fə fʌn/ & for enjoyment, not as work & để giải trí \\ \hline
\textbf{stiff} & adj & /stɪf/ & physically tight or hard to move & cứng cơ, nhức mỏi \\ \hline
\textbf{copying cartoon characters} & v phr & /ˈkɒp.i.ɪŋ kɑːˈtuːn ˈkær.ək.təz/ & drawing animated figures by looking at them & vẽ lại nhân vật hoạt hình \\ \hline
\textbf{collecting stickers} & v phr & /kəˈlektɪŋ ˈstɪk.əz/ & gathering decorative adhesive pictures & sưu tầm nhãn dán \\ \hline
\textbf{design software} & n phr & /dɪˈzaɪn ˈsɒft.weər/ & computer programs used for design & phần mềm thiết kế \\ \hline
\textbf{book a court} & v phr & /bʊk ə kɔːt/ & to reserve a sports court & đặt sân thể thao \\ \hline
\textbf{be into} & phr v & /bi ˈɪn.tuː/ & to be interested in something & thích, có hứng thú với \\ \hline
\textbf{pick up a habit} & v phr & /pɪk ʌp ə ˈhæb.ɪt/ & to gradually start doing something regularly & hình thành/thấm một thói quen \\ \hline
\end{tabularx}
\end{center}

\newpage

\section{Childhood Activities}
\begin{ieltsquestion}
\textbf{What were your favourite activities when you were a child?}
\end{ieltsquestion}
\textbf{Sample 1: Student Persona (Band 7-8)}

When I was a child, I loved playing hide-and-seek and riding my bike around the neighbourhood. Those activities were simple, but they gave me a lot of freedom and excitement.

\textbf{Sample 2: Working Professional Persona (Band 7-8)}

My favourite activities were outdoor games like football, badminton, and cycling. I spent a lot of time outside because smartphones were not as common then, so children naturally played together more.

\separator

\begin{ieltsquestion}
\textbf{What are your favourite activities?}
\end{ieltsquestion}
\textbf{Sample 1: Student Persona (Band 7-8)}

These days, I enjoy meeting friends at cafes, designing posters, and watching films. I still like going outside, but I probably spend more time on screen-based activities than I did as a child.

\textbf{Sample 2: Working Professional Persona (Band 7-8)}

Now my favourite activities are exercising, cooking, and having quiet coffee dates with friends. I prefer activities that help me recover from work rather than ones that require too much planning.

\separator

\begin{ieltsquestion}
\textbf{Did you prefer to do activities alone or with a group of people when you were a child?}
\end{ieltsquestion}
\textbf{Sample 1: Student Persona (Band 7-8)}

I preferred doing things with a group. I was quite energetic as a child, and games felt much more exciting when there were cousins or neighbours joining in.

\textbf{Sample 2: Working Professional Persona (Band 7-8)}

I definitely preferred group activities. I grew up in a neighbourhood where children played outside together every afternoon, so being alone felt a bit boring to me.

\separator

\begin{ieltsquestion}
\textbf{Are there any differences between the activities you liked when you were a child and those you like now?}
\end{ieltsquestion}
\textbf{Sample 1: Student Persona (Band 7-8)}

Yes, there are quite a few differences. As a child, I liked active outdoor games, but now I enjoy quieter activities like design, films, and cafe study sessions.

\textbf{Sample 2: Working Professional Persona (Band 7-8)}

Yes, my preferences have changed a lot. I used to love noisy group games, but now I value activities that are calming and flexible because my work life is already busy enough.

\separator


\vspace{0.8em}
\begin{center}
\renewcommand{\arraystretch}{1.65}
\rowcolors{2}{ForumFreshLight}{white}
\begin{tabularx}{\textwidth}{|V{0.22\textwidth}|C{0.07\textwidth}|I{0.20\textwidth}|Y|V{0.18\textwidth}|}
\hline
\rowcolor{ForumFreshBlue!20}
\textbf{Từ} & \textbf{Loại} & \textbf{IPA} & \textbf{English Meaning} & \textbf{Dịch nghĩa} \\
\hline
\textbf{hide-and-seek} & n & /ˌhaɪd ən ˈsiːk/ & a children's game where players hide and one person searches & trò trốn tìm \\ \hline
\textbf{neighbourhood} & n & /ˈneɪ.bə.hʊd/ & the area around where someone lives & khu phố \\ \hline
\textbf{outdoor games} & n phr & /ˈaʊt.dɔː ɡeɪmz/ & games played outside & trò chơi ngoài trời \\ \hline
\textbf{screen-based} & adj & /skriːn beɪst/ & involving phones, computers, or screens & dựa trên màn hình \\ \hline
\textbf{recover from work} & v phr & /rɪˈkʌv.ər frəm wɜːk/ & to regain energy after working & hồi sức sau công việc \\ \hline
\textbf{energetic} & adj & /ˌen.əˈdʒet.ɪk/ & full of energy & năng động \\ \hline
\textbf{joining in} & phr v & /ˌdʒɔɪ.nɪŋ ˈɪn/ & taking part in an activity & tham gia vào \\ \hline
\textbf{preference} & n & /ˈpref.ər.əns/ & something someone likes more than another & sở thích, sự ưu tiên \\ \hline
\textbf{flexible} & adj & /ˈflek.sə.bəl/ & easy to change or adapt & linh hoạt \\ \hline
\end{tabularx}
\end{center}

\newpage

\section{Sports Team}
\begin{ieltsquestion}
\textbf{Have you ever been part of a sports team?}
\end{ieltsquestion}
\textbf{Sample 1: Student Persona (Band 7-8)}

Yes, I was part of my class badminton team in high school. We were not professional or anything, but we practised after class and joined a small school tournament. It was a memorable experience because I learned how to cooperate under pressure, not just play for myself.

\textbf{Sample 2: Working Professional Persona (Band 7-8)}

Yes, I used to play on a company football team for a while. We trained once a week after work and sometimes competed with teams from other departments. I was not the best player, but it was a great way to bond with colleagues outside the office.

\separator

\begin{ieltsquestion}
\textbf{Do you like watching team games? Why?}
\end{ieltsquestion}
\textbf{Sample 1: Student Persona (Band 7-8)}

Yes, I do, especially football and volleyball matches. Team games are exciting because the result can change very quickly, and you can see how players support each other. I also enjoy the atmosphere when everyone is cheering for the same side.

\textbf{Sample 2: Working Professional Persona (Band 7-8)}

Yes, I quite enjoy watching team games, mainly because they are unpredictable. Even if one player is talented, the whole team still has to communicate well and follow a strategy. That makes the match more dramatic than just watching one person perform alone.

\separator

\begin{ieltsquestion}
\textbf{Are team sports popular in your culture?}
\end{ieltsquestion}
\textbf{Sample 1: Student Persona (Band 7-8)}

Yes, definitely. In Vietnam, football is probably the most popular team sport, and people get incredibly excited during big regional tournaments. At school, students also play basketball, volleyball, and badminton, so team sports are a normal part of student life.

\textbf{Sample 2: Working Professional Persona (Band 7-8)}

Yes, team sports are quite popular, especially football. Many companies organise football clubs or sports days because they see them as a way to build teamwork. Even people who do not play often enjoy watching national matches with friends or family.

\separator

\begin{ieltsquestion}
\textbf{What are the differences between team sports and individual sports?}
\end{ieltsquestion}
\textbf{Sample 1: Student Persona (Band 7-8)}

The biggest difference is responsibility. In team sports, you have to think about other players and adjust your actions to fit the group, while individual sports depend more on self-discipline and personal performance. I think both are useful, but they train different kinds of skills.

\textbf{Sample 2: Working Professional Persona (Band 7-8)}

Team sports require communication, trust, and coordination, whereas individual sports rely more on personal focus and consistency. In a team, one mistake can affect everyone, but one good decision can also lift the whole group. That is why I think team sports teach people a lot about collaboration.

\separator


\vspace{0.8em}
\begin{center}
\renewcommand{\arraystretch}{1.65}
\rowcolors{2}{ForumFreshLight}{white}
\begin{tabularx}{\textwidth}{|V{0.22\textwidth}|C{0.07\textwidth}|I{0.20\textwidth}|Y|V{0.18\textwidth}|}
\hline
\rowcolor{ForumFreshBlue!20}
\textbf{Từ} & \textbf{Loại} & \textbf{IPA} & \textbf{English Meaning} & \textbf{Dịch nghĩa} \\
\hline
\textbf{school tournament} & n phr & /skuːl ˈtʊə.nə.mənt/ & a sports competition at school & giải đấu ở trường \\ \hline
\textbf{cooperate} & v & /kəʊˈɒp.ə.reɪt/ & to work together with others & hợp tác \\ \hline
\textbf{bond with colleagues} & v phr & /bɒnd wɪð ˈkɒl.iːɡz/ & to build closer relationships with coworkers & gắn kết với đồng nghiệp \\ \hline
\textbf{unpredictable} & adj & /ˌʌn.prɪˈdɪk.tə.bəl/ & difficult to know what will happen & khó đoán \\ \hline
\textbf{strategy} & n & /ˈstræt.ə.dʒi/ & a planned way to achieve a goal & chiến thuật, chiến lược \\ \hline
\textbf{regional tournament} & n phr & /ˈriː.dʒən.əl ˈtʊə.nə.mənt/ & a competition within a region & giải đấu khu vực \\ \hline
\textbf{self-discipline} & n & /ˌself ˈdɪs.ə.plɪn/ & the ability to control oneself and keep working & tính kỷ luật tự thân \\ \hline
\textbf{coordination} & n & /kəʊˌɔː.dɪˈneɪ.ʃən/ & the ability to organise actions together & sự phối hợp \\ \hline
\textbf{collaboration} & n & /kəˌlæb.əˈreɪ.ʃən/ & working together to achieve something & sự cộng tác \\ \hline
\end{tabularx}
\end{center}

\newpage

\section{Reading}
\begin{ieltsquestion}
\textbf{Do you like reading?}
\end{ieltsquestion}
\textbf{Sample 1: Student Persona (Band 7-8)}

Yes, but my reading habits have changed quite a lot. I do not read novels as much as I used to, but I read articles, course materials, and short posts related to business or design almost every day. Reading helps me pick up new ideas without feeling like I am forcing myself to study.

\textbf{Sample 2: Working Professional Persona (Band 7-8)}

Yes, I do, although I have less time for it now. I mostly read business books, newsletters, and articles about marketing or productivity. Even twenty minutes of reading before bed helps me slow down and learn something useful.

\separator

\begin{ieltsquestion}
\textbf{Do you prefer to read on paper or on a screen?}
\end{ieltsquestion}
\textbf{Sample 1: Student Persona (Band 7-8)}

For textbooks, I prefer reading on paper because it is easier to highlight important points and write notes in the margin. However, for quick articles or lecture slides, reading on a screen is more convenient. So I would say it depends on the purpose.

\textbf{Sample 2: Working Professional Persona (Band 7-8)}

I prefer screens for work because documents are easier to search, edit, and share. But when I read for pleasure, I still like printed books because they feel less distracting. A physical book also helps me stay away from notifications for a while.

\separator

\begin{ieltsquestion}
\textbf{Do you prefer scanning or detailed reading?}
\end{ieltsquestion}
\textbf{Sample 1: Student Persona (Band 7-8)}

It depends on what I am reading. If I only need the main idea, I scan the text first to save time, especially before class discussions. But for exam preparation, I have to read carefully because small details can make a big difference.

\textbf{Sample 2: Working Professional Persona (Band 7-8)}

At work, I often scan reports first because I need to find key numbers or conclusions quickly. After that, if something looks important, I go back and read it in detail. This approach helps me avoid wasting time on information that is not relevant.

\separator

\begin{ieltsquestion}
\textbf{When do you need to read carefully, and when not?}
\end{ieltsquestion}
\textbf{Sample 1: Student Persona (Band 7-8)}

I need to read carefully when I study academic materials, exam instructions, or assignment requirements. If I miss one small condition, I might lose marks. But when I read social media posts or simple news updates, scanning is usually enough.

\textbf{Sample 2: Working Professional Persona (Band 7-8)}

I read carefully when dealing with contracts, client briefs, or important emails because misunderstanding one sentence could cause real problems. On the other hand, I do not read every newsletter or casual update word by word. For those, I just skim for anything useful.

\separator


\vspace{0.8em}
\begin{center}
\renewcommand{\arraystretch}{1.65}
\rowcolors{2}{ForumFreshLight}{white}
\begin{tabularx}{\textwidth}{|V{0.22\textwidth}|C{0.07\textwidth}|I{0.20\textwidth}|Y|V{0.18\textwidth}|}
\hline
\rowcolor{ForumFreshBlue!20}
\textbf{Từ} & \textbf{Loại} & \textbf{IPA} & \textbf{English Meaning} & \textbf{Dịch nghĩa} \\
\hline
\textbf{reading habit} & n phr & /ˈriː.dɪŋ ˈhæb.ɪt/ & a regular way of reading & thói quen đọc \\ \hline
\textbf{course materials} & n phr & /kɔːs məˈtɪə.ri.əlz/ & learning resources for a subject & tài liệu học tập \\ \hline
\textbf{newsletter} & n & /ˈnjuːzˌlet.ər/ & an email or document with updates & bản tin định kỳ \\ \hline
\textbf{highlight} & v & /ˈhaɪ.laɪt/ & to mark important information & đánh dấu nổi bật \\ \hline
\textbf{margin} & n & /ˈmɑː.dʒɪn/ & the blank space at the side of a page & lề trang \\ \hline
\textbf{printed book} & n phr & /ˈprɪn.tɪd bʊk/ & a physical book on paper & sách in \\ \hline
\textbf{scan the text} & v phr & /skæn ðə tekst/ & to read quickly for key information & đọc lướt để tìm ý chính \\ \hline
\textbf{exam preparation} & n phr & /ɪɡˈzæm ˌprep.əˈreɪ.ʃən/ & studying before an exam & ôn thi \\ \hline
\textbf{client brief} & n phr & /ˈklaɪ.ənt briːf/ & instructions or requirements from a client & bản yêu cầu của khách hàng \\ \hline
\textbf{skim} & v & /skɪm/ & to read quickly without focusing on every word & đọc lướt \\ \hline
\end{tabularx}
\end{center}

\newpage

\section{Typing}
\begin{ieltsquestion}
\textbf{Do you prefer typing or handwriting?}
\end{ieltsquestion}
\textbf{Sample 1: Student Persona (Band 7-8)}

I prefer typing for most school tasks because it is much faster and easier to edit. When I write essays or prepare group reports, typing helps me organise ideas more clearly. That said, I still use handwriting when I need to memorise something because it feels more personal and focused.

\textbf{Sample 2: Working Professional Persona (Band 7-8)}

I definitely prefer typing because almost everything at work is digital now. Emails, reports, meeting notes, and project updates all need to be typed and shared quickly. Handwriting is still useful for quick brainstorming, but it is not practical for my daily workload.

\separator

\begin{ieltsquestion}
\textbf{When did you learn how to type on a keyboard?}
\end{ieltsquestion}
\textbf{Sample 1: Student Persona (Band 7-8)}

I think I learned properly when I was in secondary school. At first, I typed very slowly and looked at the keyboard all the time, but I improved naturally by doing homework and chatting with friends online.

\textbf{Sample 2: Working Professional Persona (Band 7-8)}

I started learning in high school, but I became much faster at university when I had to write assignments regularly. Later, work made my typing even more efficient because I had to reply to emails and prepare documents under time pressure.

\separator

\begin{ieltsquestion}
\textbf{Do you type on a desktop or laptop keyboard every day?}
\end{ieltsquestion}
\textbf{Sample 1: Student Persona (Band 7-8)}

Yes, I type on my laptop almost every day. I use it for taking notes, writing assignments, making presentations, and sometimes doing freelance design work. Without a keyboard, my study routine would be much less efficient.

\textbf{Sample 2: Working Professional Persona (Band 7-8)}

Yes, I use a laptop keyboard every working day, and sometimes an external keyboard at the office. I type so much that I have become quite sensitive to whether a keyboard feels comfortable or not.

\separator

\begin{ieltsquestion}
\textbf{How do you improve your typing?}
\end{ieltsquestion}
\textbf{Sample 1: Student Persona (Band 7-8)}

I improve mostly through practice rather than formal training. The more I type essays, messages, and notes, the faster I become. I also try to type without looking at the keyboard because that saves a surprising amount of time.

\textbf{Sample 2: Working Professional Persona (Band 7-8)}

For me, accuracy matters as much as speed, so I try to slow down slightly when writing important emails. I also use keyboard shortcuts and templates for repetitive tasks. Over time, those small habits make typing feel smoother and more professional.

\separator


\vspace{0.8em}
\begin{center}
\renewcommand{\arraystretch}{1.65}
\rowcolors{2}{ForumFreshLight}{white}
\begin{tabularx}{\textwidth}{|V{0.22\textwidth}|C{0.07\textwidth}|I{0.20\textwidth}|Y|V{0.18\textwidth}|}
\hline
\rowcolor{ForumFreshBlue!20}
\textbf{Từ} & \textbf{Loại} & \textbf{IPA} & \textbf{English Meaning} & \textbf{Dịch nghĩa} \\
\hline
\textbf{edit} & v & /ˈed.ɪt/ & to change or correct written content & chỉnh sửa \\ \hline
\textbf{memorise} & v & /ˈmem.ə.raɪz/ & to learn something so you can remember it & ghi nhớ \\ \hline
\textbf{brainstorming} & n & /ˈbreɪnˌstɔː.mɪŋ/ & quickly producing ideas & động não, lên ý tưởng \\ \hline
\textbf{practical} & adj & /ˈpræk.tɪ.kəl/ & useful and suitable for real situations & thực tế, tiện dụng \\ \hline
\textbf{under time pressure} & adv phr & /ˈʌn.də taɪm ˈpreʃ.ər/ & with limited time available & dưới áp lực thời gian \\ \hline
\textbf{external keyboard} & n phr & /ɪkˈstɜː.nəl ˈkiː.bɔːd/ & a separate keyboard connected to a device & bàn phím rời \\ \hline
\textbf{accuracy} & n & /ˈæk.jə.rə.si/ & correctness without mistakes & độ chính xác \\ \hline
\textbf{keyboard shortcut} & n phr & /ˈkiː.bɔːd ˈʃɔːt.kʌt/ & a key combination that performs an action quickly & phím tắt \\ \hline
\textbf{repetitive task} & n phr & /rɪˈpet.ə.tɪv tɑːsk/ & a task done again and again & nhiệm vụ lặp đi lặp lại \\ \hline
\end{tabularx}
\end{center}

\newpage

\section*{PART 1: VIEWS, PLACES \& TRAVEL}
\addcontentsline{toc}{section}{PART 1: VIEWS, PLACES \& TRAVEL}

\section{Views}
\begin{ieltsquestion}
\textbf{Have you seen an unforgettable and beautiful view or scenery?}
\end{ieltsquestion}
\textbf{Sample 1: Student Persona (Band 7-8)}

Yes, I once watched the sunrise from My Khe Beach in Da Nang, and it was genuinely unforgettable. The sky slowly turned pink and orange, and the sea looked almost silver in the early light. It was a simple moment, but it made me feel unusually calm before a busy school day.

\textbf{Sample 2: Working Professional Persona (Band 7-8)}

Yes, the most memorable view I have seen was from a rooftop bar in Ho Chi Minh City at night. The whole skyline was lit up, and you could see streams of traffic moving below. It made the city look chaotic but strangely beautiful at the same time.

\separator

\begin{ieltsquestion}
\textbf{Do you like taking pictures of different views?}
\end{ieltsquestion}
\textbf{Sample 1: Student Persona (Band 7-8)}

Yes, I do, especially when I travel or go to the beach with friends. I am not a professional photographer, but I like capturing views because photos help me remember how I felt at that moment. Sometimes I also use those pictures as design inspiration.

\textbf{Sample 2: Working Professional Persona (Band 7-8)}

Yes, but I do it more casually now. If I see a nice sunset, a quiet street, or an interesting city view, I take a quick photo on my phone. I rarely post everything online, but I like keeping those images as small memories.

\separator

\begin{ieltsquestion}
\textbf{Do you prefer views in urban areas or rural areas?}
\end{ieltsquestion}
\textbf{Sample 1: Student Persona (Band 7-8)}

I prefer rural or natural views because they feel more relaxing. Urban views can be impressive, especially at night, but nature gives me a stronger sense of peace. When I am stressed with exams, seeing the sea or mountains helps me clear my mind.

\textbf{Sample 2: Working Professional Persona (Band 7-8)}

I enjoy both, but I slightly prefer rural views because my daily life is already surrounded by buildings and traffic. A view of rice fields, hills, or a quiet lake feels like a real escape. It reminds me that life does not always have to move so fast.

\separator

\begin{ieltsquestion}
\textbf{Do you prefer views in your own country or in other countries?}
\end{ieltsquestion}
\textbf{Sample 1: Student Persona (Band 7-8)}

For now, I prefer views in my own country because I have not travelled abroad much. Vietnam has beaches, mountains, old towns, and rice terraces, so there is already a lot to explore. I would love to see foreign landscapes too, but I do not feel I have fully discovered my own country yet.

\textbf{Sample 2: Working Professional Persona (Band 7-8)}

I appreciate both, but views in my own country feel more meaningful to me. When I see places like Ha Long Bay or the mountains in Da Lat, there is a sense of familiarity and pride. Foreign views can be stunning, but local scenery often feels closer to my heart.

\separator


\vspace{0.8em}
\begin{center}
\renewcommand{\arraystretch}{1.65}
\rowcolors{2}{ForumFreshLight}{white}
\begin{tabularx}{\textwidth}{|V{0.22\textwidth}|C{0.07\textwidth}|I{0.20\textwidth}|Y|V{0.18\textwidth}|}
\hline
\rowcolor{ForumFreshBlue!20}
\textbf{Từ} & \textbf{Loại} & \textbf{IPA} & \textbf{English Meaning} & \textbf{Dịch nghĩa} \\
\hline
\textbf{unforgettable} & adj & /ˌʌn.fəˈɡet.ə.bəl/ & impossible to forget & không thể quên \\ \hline
\textbf{sunrise} & n & /ˈsʌn.raɪz/ & the time when the sun appears in the morning & bình minh \\ \hline
\textbf{skyline} & n & /ˈskaɪ.laɪn/ & the outline of buildings against the sky & đường chân trời đô thị \\ \hline
\textbf{stream of traffic} & n phr & /striːm əv ˈtræf.ɪk/ & a continuous flow of vehicles & dòng xe cộ \\ \hline
\textbf{capture} & v & /ˈkæp.tʃər/ & to record something in a photo or video & ghi lại, chụp lại \\ \hline
\textbf{design inspiration} & n phr & /dɪˈzaɪn ˌɪn.spəˈreɪ.ʃən/ & ideas that help creative design & cảm hứng thiết kế \\ \hline
\textbf{sense of peace} & n phr & /sens əv piːs/ & a feeling of calmness & cảm giác bình yên \\ \hline
\textbf{rice terrace} & n phr & /raɪs ˈter.əs/ & stepped fields used for growing rice & ruộng bậc thang \\ \hline
\textbf{familiarity} & n & /fəˌmɪl.iˈær.ə.ti/ & the feeling of knowing something well & sự thân thuộc \\ \hline
\textbf{stunning} & adj & /ˈstʌn.ɪŋ/ & extremely beautiful or impressive & tuyệt đẹp, ấn tượng \\ \hline
\end{tabularx}
\end{center}

\newpage

\section{Scenery}
\begin{ieltsquestion}
\textbf{Do you look out the window at the scenery when travelling by bus or car?}
\end{ieltsquestion}
\textbf{Sample 1: Student Persona (Band 7-8)}

Yes, I do it almost automatically, especially when I travel by bus between my hometown and university. Looking at the roads, fields, and small houses outside makes the journey feel shorter. It is also a nice break from staring at my phone for hours.

\textbf{Sample 2: Working Professional Persona (Band 7-8)}

Yes, I often look out the window, particularly when I am not driving myself. Even in the city, I like observing little details such as street vendors, old buildings, and changing neighbourhoods. It gives me something calm to focus on during traffic.

\separator

\begin{ieltsquestion}
\textbf{What are the most beautiful sights you have seen while travelling?}
\end{ieltsquestion}
\textbf{Sample 1: Student Persona (Band 7-8)}

One of the most beautiful sights I have seen was the view from Hai Van Pass. You can see the mountains on one side and the sea on the other, which feels incredibly dramatic. I remember taking lots of photos because the scenery looked different every few minutes.

\textbf{Sample 2: Working Professional Persona (Band 7-8)}

For me, the most beautiful sight was in Da Lat, where the hills were covered in mist early in the morning. It was quiet, cool, and completely different from the busy city where I work. That kind of scenery makes me feel as if time slows down.

\separator

\begin{ieltsquestion}
\textbf{Do you prefer the mountains or the sea?}
\end{ieltsquestion}
\textbf{Sample 1: Student Persona (Band 7-8)}

I slightly prefer the sea because I grew up near the coast. The sound of waves and the open horizon make me feel relaxed, especially when I am stressed about school. Mountains are beautiful too, but the sea feels more familiar to me.

\textbf{Sample 2: Working Professional Persona (Band 7-8)}

I prefer the mountains these days. Since I spend most of my time in a hot, crowded city, cool mountain air feels refreshing. I also like the slower pace in mountain towns, where people seem less rushed.

\separator

\begin{ieltsquestion}
\textbf{Do you like to take scenery pictures?}
\end{ieltsquestion}
\textbf{Sample 1: Student Persona (Band 7-8)}

Yes, I like taking scenery pictures, although I mostly use my phone. I enjoy capturing sunsets, beaches, and quiet streets because those photos can bring back the mood of a trip. Sometimes I even use them as inspiration for my design projects.

\textbf{Sample 2: Working Professional Persona (Band 7-8)}

Yes, but I do not take photos obsessively. If a view is really striking, I take a few pictures and then put my phone away so I can enjoy it properly. I think photos are nice, but they should not replace the actual experience.

\separator


\vspace{0.8em}
\begin{center}
\renewcommand{\arraystretch}{1.65}
\rowcolors{2}{ForumFreshLight}{white}
\begin{tabularx}{\textwidth}{|V{0.22\textwidth}|C{0.07\textwidth}|I{0.20\textwidth}|Y|V{0.18\textwidth}|}
\hline
\rowcolor{ForumFreshBlue!20}
\textbf{Từ} & \textbf{Loại} & \textbf{IPA} & \textbf{English Meaning} & \textbf{Dịch nghĩa} \\
\hline
\textbf{automatically} & adv & /ˌɔː.təˈmæt.ɪ.kəl.i/ & without thinking much & một cách tự nhiên, tự động \\ \hline
\textbf{street vendor} & n phr & /striːt ˈven.dər/ & someone who sells goods on the street & người bán hàng rong \\ \hline
\textbf{dramatic} & adj & /drəˈmæt.ɪk/ & impressive and exciting to look at & ấn tượng, hùng vĩ \\ \hline
\textbf{mist} & n & /mɪst/ & thin cloud close to the ground & sương mù mỏng \\ \hline
\textbf{open horizon} & n phr & /ˈəʊ.pən həˈraɪ.zən/ & a wide view where the sky meets the land or sea & đường chân trời rộng mở \\ \hline
\textbf{refreshing} & adj & /rɪˈfreʃ.ɪŋ/ & making someone feel fresh and energetic & sảng khoái \\ \hline
\textbf{capture} & v & /ˈkæp.tʃər/ & to record something in a photo & ghi lại, chụp lại \\ \hline
\textbf{striking} & adj & /ˈstraɪ.kɪŋ/ & very noticeable or impressive & nổi bật, gây ấn tượng mạnh \\ \hline
\textbf{actual experience} & n phr & /ˈæk.tʃu.əl ɪkˈspɪə.ri.əns/ & the real moment or event itself & trải nghiệm thực tế \\ \hline
\end{tabularx}
\end{center}

\newpage

\section{Building}
\begin{ieltsquestion}
\textbf{Are there tall buildings near your home?}
\end{ieltsquestion}
\textbf{Sample 1: Student Persona (Band 7-8)}

There are a few tall apartment buildings near where I live, but the area is not packed with skyscrapers. Most buildings around my university are medium-sized, like cafes, dormitories, and small offices. I actually prefer that because the streets feel more open and less overwhelming.

\textbf{Sample 2: Working Professional Persona (Band 7-8)}

Yes, there are quite a few tall buildings near my apartment. I live close to a business district, so there are office towers, hotels, and modern apartment blocks all around. It looks impressive at night, but it also means the area can be crowded during rush hour.

\separator

\begin{ieltsquestion}
\textbf{Do you take photos of buildings?}
\end{ieltsquestion}
\textbf{Sample 1: Student Persona (Band 7-8)}

Sometimes, especially if the building has an interesting design or colourful walls. As someone who likes graphic design, I notice shapes, patterns, and lighting quite a lot. I do not post every picture, but I keep them for inspiration.

\textbf{Sample 2: Working Professional Persona (Band 7-8)}

Yes, occasionally. When I travel, I like taking photos of old buildings, temples, or unusual modern architecture. They tell you a lot about a place's history and personality, not just how it looks.

\separator

\begin{ieltsquestion}
\textbf{Is there a building that you would like to visit?}
\end{ieltsquestion}
\textbf{Sample 1: Student Persona (Band 7-8)}

I would like to visit the Bitexco Financial Tower in Ho Chi Minh City. I have seen it in photos many times, and I think the observation deck would give me a great view of the city. It would also be interesting to compare that skyline with the quieter cityscape in Da Nang.

\textbf{Sample 2: Working Professional Persona (Band 7-8)}

I would love to visit the Sydney Opera House one day. It is not just famous; its design is so distinctive that you can recognise it immediately. I think seeing it in person would be much more impressive than seeing it online.

\separator

\begin{ieltsquestion}
\textbf{Do you want to live in a tall building?}
\end{ieltsquestion}
\textbf{Sample 1: Student Persona (Band 7-8)}

Maybe in the future, but not necessarily now. Living in a tall building could be convenient and offer a nice view, but I would worry about high rent and crowded lifts. As a student, practicality matters more than having a fancy address.

\textbf{Sample 2: Working Professional Persona (Band 7-8)}

Yes, I would not mind living in a tall building if it were well-managed. A high-floor apartment can offer better views, more security, and useful facilities like a gym or swimming pool. However, I would still want the building to be quiet and not too densely populated.

\separator


\vspace{0.8em}
\begin{center}
\renewcommand{\arraystretch}{1.65}
\rowcolors{2}{ForumFreshLight}{white}
\begin{tabularx}{\textwidth}{|V{0.22\textwidth}|C{0.07\textwidth}|I{0.20\textwidth}|Y|V{0.18\textwidth}|}
\hline
\rowcolor{ForumFreshBlue!20}
\textbf{Từ} & \textbf{Loại} & \textbf{IPA} & \textbf{English Meaning} & \textbf{Dịch nghĩa} \\
\hline
\textbf{skyscraper} & n & /ˈskaɪˌskreɪ.pər/ & a very tall building & nhà chọc trời \\ \hline
\textbf{overwhelming} & adj & /ˌəʊ.vəˈwel.mɪŋ/ & too intense or crowded to handle easily & choáng ngợp \\ \hline
\textbf{business district} & n phr & /ˈbɪz.nəs ˈdɪs.trɪkt/ & an area with many offices and companies & khu thương mại/văn phòng \\ \hline
\textbf{apartment block} & n phr & /əˈpɑːt.mənt blɒk/ & a large building with many apartments & tòa chung cư \\ \hline
\textbf{pattern} & n & /ˈpæt.ən/ & a repeated design or arrangement & họa tiết, kiểu mẫu \\ \hline
\textbf{architecture} & n & /ˈɑː.kɪ.tek.tʃər/ & the design and style of buildings & kiến trúc \\ \hline
\textbf{observation deck} & n phr & /ˌɒb.zəˈveɪ.ʃən dek/ & a high area for viewing a city or landscape & đài quan sát \\ \hline
\textbf{cityscape} & n & /ˈsɪt.i.skeɪp/ & the visual appearance of a city & cảnh quan thành phố \\ \hline
\textbf{distinctive} & adj & /dɪˈstɪŋk.tɪv/ & clearly different and easy to recognise & đặc trưng, dễ nhận ra \\ \hline
\textbf{densely populated} & adj phr & /ˈdens.li ˈpɒp.jə.leɪ.tɪd/ & having many people in a small area & đông dân, mật độ cao \\ \hline
\end{tabularx}
\end{center}

\newpage

\section{Museum}
\begin{ieltsquestion}
\textbf{Do you think museums are important?}
\end{ieltsquestion}
\textbf{Sample 1: Student Persona (Band 7-8)}

Yes, I think museums are important because they make history and culture easier to understand. Reading about the past in a textbook can feel quite dry, but seeing real objects, photos, and old documents makes the story more vivid. For students like me, museums can turn learning into something more memorable.

\textbf{Sample 2: Working Professional Persona (Band 7-8)}

Yes, definitely. Museums preserve things that might otherwise be forgotten, such as traditional crafts, war memories, or local customs. I also think they give busy adults a chance to slow down and reconnect with history instead of only focusing on work and modern life.

\separator

\begin{ieltsquestion}
\textbf{Are there many museums in your hometown?}
\end{ieltsquestion}
\textbf{Sample 1: Student Persona (Band 7-8)}

There are a few, but I would not say there are many. In Da Nang, for example, there are museums about Cham sculpture, local history, and fine arts. They are not huge, but they are still worth visiting if someone wants to understand the city better.

\textbf{Sample 2: Working Professional Persona (Band 7-8)}

My hometown has some interesting museums, though not as many as a capital city. The most famous ones are connected to local history and culture. I think they could attract more visitors if they became more interactive and better promoted online.

\separator

\begin{ieltsquestion}
\textbf{Do you often visit a museum?}
\end{ieltsquestion}
\textbf{Sample 1: Student Persona (Band 7-8)}

Not very often, to be honest. I usually visit museums when my teachers organise a field trip or when friends from another city come to visit. I enjoy the experience once I am there, but it is not something I do every weekend.

\textbf{Sample 2: Working Professional Persona (Band 7-8)}

I do not visit museums regularly, mainly because my schedule is quite full. However, when I travel, I try to include at least one museum because it helps me understand the place beyond restaurants and tourist spots.

\separator

\begin{ieltsquestion}
\textbf{When was the last time you visited a museum?}
\end{ieltsquestion}
\textbf{Sample 1: Student Persona (Band 7-8)}

The last time was a few months ago when I visited the Museum of Cham Sculpture with some classmates. We had to do a small assignment about local culture, but I ended up enjoying it more than expected. The carvings were detailed, and the atmosphere was very quiet.

\textbf{Sample 2: Working Professional Persona (Band 7-8)}

The last time was during a short trip to Hanoi last year. I visited the Vietnam Museum of Ethnology, and I was impressed by how it showed the lifestyles of different ethnic groups. It was educational without feeling too heavy.

\separator


\vspace{0.8em}
\begin{center}
\renewcommand{\arraystretch}{1.65}
\rowcolors{2}{ForumFreshLight}{white}
\begin{tabularx}{\textwidth}{|V{0.22\textwidth}|C{0.07\textwidth}|I{0.20\textwidth}|Y|V{0.18\textwidth}|}
\hline
\rowcolor{ForumFreshBlue!20}
\textbf{Từ} & \textbf{Loại} & \textbf{IPA} & \textbf{English Meaning} & \textbf{Dịch nghĩa} \\
\hline
\textbf{preserve} & v & /prɪˈzɜːv/ & to protect something from being lost or damaged & bảo tồn \\ \hline
\textbf{vivid} & adj & /ˈvɪv.ɪd/ & clear, detailed, and easy to imagine & sinh động \\ \hline
\textbf{reconnect with history} & v phr & /ˌriː.kəˈnekt wɪð ˈhɪs.tər.i/ & to feel connected to the past again & kết nối lại với lịch sử \\ \hline
\textbf{Cham sculpture} & n phr & /tʃɑːm ˈskʌlp.tʃər/ & carved art from Cham culture & điêu khắc Chăm \\ \hline
\textbf{interactive} & adj & /ˌɪn.təˈræk.tɪv/ & involving active participation & có tính tương tác \\ \hline
\textbf{field trip} & n phr & /fiːld trɪp/ & an educational visit outside school & chuyến đi thực tế \\ \hline
\textbf{tourist spot} & n phr & /ˈtʊə.rɪst spɒt/ & a place popular with visitors & điểm du lịch \\ \hline
\textbf{carving} & n & /ˈkɑː.vɪŋ/ & an object or pattern cut from wood or stone & tác phẩm chạm khắc \\ \hline
\textbf{ethnic group} & n phr & /ˈeθ.nɪk ɡruːp/ & a group sharing culture or ancestry & nhóm dân tộc \\ \hline
\end{tabularx}
\end{center}

\newpage

\section{Public Places}
\begin{ieltsquestion}
\textbf{Do you often go to public places with your friends?}
\end{ieltsquestion}
\textbf{Sample 1: Student Persona (Band 7-8)}

Yes, quite often. My friends and I usually go to cafes, parks, shopping centres, or the beach because those places are affordable and easy to reach. Public places are also more comfortable for students because no one has to host everyone at home.

\textbf{Sample 2: Working Professional Persona (Band 7-8)}

Yes, but mostly on weekends. My friends and I often meet at coffee shops, bookstores, or walking streets because they are convenient meeting points. After a long workweek, I prefer public places that are lively but not too noisy.

\separator

\begin{ieltsquestion}
\textbf{Have you ever talked with someone you don't know in public places?}
\end{ieltsquestion}
\textbf{Sample 1: Student Persona (Band 7-8)}

Yes, a few times. For example, I once talked to a tourist at the beach because she needed directions, and we ended up chatting briefly about local food. It was not a deep conversation, but it felt friendly and natural.

\textbf{Sample 2: Working Professional Persona (Band 7-8)}

Yes, occasionally. I have spoken to strangers in cafes, airports, and queues, usually about simple things like directions or whether a seat is taken. I do not start conversations all the time, but I do not mind them when they happen naturally.

\separator

\begin{ieltsquestion}
\textbf{Do you wear headphones in public places?}
\end{ieltsquestion}
\textbf{Sample 1: Student Persona (Band 7-8)}

Yes, I often wear headphones when I am on the bus or studying in a cafe. They help me block out background noise and focus better. But if I am walking on a busy street, I keep the volume low for safety.

\textbf{Sample 2: Working Professional Persona (Band 7-8)}

Yes, especially when I commute or work in a shared space. Headphones help me create a small private bubble in a crowded environment. Still, I try to stay aware of what is happening around me, especially in traffic.

\separator

\begin{ieltsquestion}
\textbf{Would you like to see more public places near where you live?}
\end{ieltsquestion}
\textbf{Sample 1: Student Persona (Band 7-8)}

Yes, I would love to see more green spaces and study-friendly public areas. Many students need places where they can relax or work without spending too much money. A clean park or public reading space would be really useful.

\textbf{Sample 2: Working Professional Persona (Band 7-8)}

Yes, definitely. I think cities need more parks, benches, and pedestrian-friendly areas, not just malls and office buildings. Good public spaces make daily life healthier and give people somewhere to breathe outside their apartments.

\separator


\vspace{0.8em}
\begin{center}
\renewcommand{\arraystretch}{1.65}
\rowcolors{2}{ForumFreshLight}{white}
\begin{tabularx}{\textwidth}{|V{0.22\textwidth}|C{0.07\textwidth}|I{0.20\textwidth}|Y|V{0.18\textwidth}|}
\hline
\rowcolor{ForumFreshBlue!20}
\textbf{Từ} & \textbf{Loại} & \textbf{IPA} & \textbf{English Meaning} & \textbf{Dịch nghĩa} \\
\hline
\textbf{affordable} & adj & /əˈfɔː.də.bəl/ & reasonably priced & giá cả phải chăng \\ \hline
\textbf{host} & v & /həʊst/ & to invite and receive guests & tiếp đón, tổ chức tại nhà \\ \hline
\textbf{meeting point} & n phr & /ˈmiː.tɪŋ pɔɪnt/ & a place where people agree to meet & điểm hẹn \\ \hline
\textbf{directions} & n & /daɪˈrek.ʃənz/ & instructions for how to get somewhere & chỉ đường \\ \hline
\textbf{queue} & n & /kjuː/ & a line of people waiting & hàng chờ \\ \hline
\textbf{block out} & phr v & /blɒk aʊt/ & to stop something from being heard or noticed & chặn, loại bỏ \\ \hline
\textbf{private bubble} & n phr & /ˈpraɪ.vət ˈbʌb.əl/ & a personal mental space in public & không gian riêng tư nhỏ \\ \hline
\textbf{green space} & n phr & /ɡriːn speɪs/ & an area with grass, trees, or plants & không gian xanh \\ \hline
\textbf{pedestrian-friendly} & adj & /pəˈdes.tri.ən ˈfrend.li/ & safe and pleasant for walkers & thân thiện với người đi bộ \\ \hline
\textbf{somewhere to breathe} & n phr & /ˈsʌm.weə tuː briːð/ & a place to relax and feel less crowded & nơi để thư giãn, dễ thở hơn \\ \hline
\end{tabularx}
\end{center}

\newpage

\section{Crowded Place}
\begin{ieltsquestion}
\textbf{Is the city where you live crowded?}
\end{ieltsquestion}
\textbf{Sample 1: Student Persona (Band 7-8)}

[Nhập câu trả lời vào đây]

\textbf{Sample 2: Working Professional Persona (Band 7-8)}

[Nhập câu trả lời vào đây]

\separator

\begin{ieltsquestion}
\textbf{Is there a crowded place near where you live?}
\end{ieltsquestion}
\textbf{Sample 1: Student Persona (Band 7-8)}

[Nhập câu trả lời vào đây]

\textbf{Sample 2: Working Professional Persona (Band 7-8)}

[Nhập câu trả lời vào đây]

\separator

\begin{ieltsquestion}
\textbf{Do you like crowded places?}
\end{ieltsquestion}
\textbf{Sample 1: Student Persona (Band 7-8)}

[Nhập câu trả lời vào đây]

\textbf{Sample 2: Working Professional Persona (Band 7-8)}

[Nhập câu trả lời vào đây]

\separator

\begin{ieltsquestion}
\textbf{When was the last time you were in a crowded place?}
\end{ieltsquestion}
\textbf{Sample 1: Student Persona (Band 7-8)}

[Nhập câu trả lời vào đây]

\textbf{Sample 2: Working Professional Persona (Band 7-8)}

[Nhập câu trả lời vào đây]

\separator

\begin{ieltsquestion}
\textbf{Do most people like crowded places?}
\end{ieltsquestion}
\textbf{Sample 1: Student Persona (Band 7-8)}

[Nhập câu trả lời vào đây]

\textbf{Sample 2: Working Professional Persona (Band 7-8)}

[Nhập câu trả lời vào đây]

\separator


\vspace{0.8em}
\begin{center}
\renewcommand{\arraystretch}{1.65}
\rowcolors{2}{ForumFreshLight}{white}
\begin{tabularx}{\textwidth}{|V{0.22\textwidth}|C{0.07\textwidth}|I{0.20\textwidth}|Y|V{0.18\textwidth}|}
\hline
\rowcolor{ForumFreshBlue!20}
\textbf{Từ} & \textbf{Loại} & \textbf{IPA} & \textbf{English Meaning} & \textbf{Dịch nghĩa} \\
\hline
\textbf{word} & & /pron/ & & Nghĩa \\ \hline
\end{tabularx}
\end{center}

\newpage

\section*{PART 1: SOCIAL \& EMOTIONAL TOPICS}
\addcontentsline{toc}{section}{PART 1: SOCIAL \& EMOTIONAL TOPICS}

\section{Chatting / Conversation}
\begin{ieltsquestion}
\textbf{Do you like chatting with friends?}
\end{ieltsquestion}
\textbf{Sample 1: Student Persona (Band 7-8)}

[Nhập câu trả lời vào đây]

\textbf{Sample 2: Working Professional Persona (Band 7-8)}

[Nhập câu trả lời vào đây]

\separator

\begin{ieltsquestion}
\textbf{What do you usually chat about with friends?}
\end{ieltsquestion}
\textbf{Sample 1: Student Persona (Band 7-8)}

[Nhập câu trả lời vào đây]

\textbf{Sample 2: Working Professional Persona (Band 7-8)}

[Nhập câu trả lời vào đây]

\separator

\begin{ieltsquestion}
\textbf{Do you prefer to chat with a group of people or with only one friend?}
\end{ieltsquestion}
\textbf{Sample 1: Student Persona (Band 7-8)}

[Nhập câu trả lời vào đây]

\textbf{Sample 2: Working Professional Persona (Band 7-8)}

[Nhập câu trả lời vào đây]

\separator

\begin{ieltsquestion}
\textbf{Do you prefer to communicate face-to-face or via social media?}
\end{ieltsquestion}
\textbf{Sample 1: Student Persona (Band 7-8)}

[Nhập câu trả lời vào đây]

\textbf{Sample 2: Working Professional Persona (Band 7-8)}

[Nhập câu trả lời vào đây]

\separator

\begin{ieltsquestion}
\textbf{Do you argue with friends?}
\end{ieltsquestion}
\textbf{Sample 1: Student Persona (Band 7-8)}

[Nhập câu trả lời vào đây]

\textbf{Sample 2: Working Professional Persona (Band 7-8)}

[Nhập câu trả lời vào đây]

\separator


\vspace{0.8em}
\begin{center}
\renewcommand{\arraystretch}{1.65}
\rowcolors{2}{ForumFreshLight}{white}
\begin{tabularx}{\textwidth}{|V{0.22\textwidth}|C{0.07\textwidth}|I{0.20\textwidth}|Y|V{0.18\textwidth}|}
\hline
\rowcolor{ForumFreshBlue!20}
\textbf{Từ} & \textbf{Loại} & \textbf{IPA} & \textbf{English Meaning} & \textbf{Dịch nghĩa} \\
\hline
\textbf{word} & & /pron/ & & Nghĩa \\ \hline
\end{tabularx}
\end{center}

\newpage

\section{Sharing}
\begin{ieltsquestion}
\textbf{Did your parents teach you to share when you were a child?}
\end{ieltsquestion}
\textbf{Sample 1: Student Persona (Band 7-8)}

[Nhập câu trả lời vào đây]

\textbf{Sample 2: Working Professional Persona (Band 7-8)}

[Nhập câu trả lời vào đây]

\separator

\begin{ieltsquestion}
\textbf{What kind of things do you like to share with others?}
\end{ieltsquestion}
\textbf{Sample 1: Student Persona (Band 7-8)}

[Nhập câu trả lời vào đây]

\textbf{Sample 2: Working Professional Persona (Band 7-8)}

[Nhập câu trả lời vào đây]

\separator

\begin{ieltsquestion}
\textbf{What kind of things are not suitable for sharing?}
\end{ieltsquestion}
\textbf{Sample 1: Student Persona (Band 7-8)}

[Nhập câu trả lời vào đây]

\textbf{Sample 2: Working Professional Persona (Band 7-8)}

[Nhập câu trả lời vào đây]

\separator

\begin{ieltsquestion}
\textbf{Do you have anything to share with others recently?}
\end{ieltsquestion}
\textbf{Sample 1: Student Persona (Band 7-8)}

[Nhập câu trả lời vào đây]

\textbf{Sample 2: Working Professional Persona (Band 7-8)}

[Nhập câu trả lời vào đây]

\separator

\begin{ieltsquestion}
\textbf{Who is the first person you would like to share good news with?}
\end{ieltsquestion}
\textbf{Sample 1: Student Persona (Band 7-8)}

[Nhập câu trả lời vào đây]

\textbf{Sample 2: Working Professional Persona (Band 7-8)}

[Nhập câu trả lời vào đây]

\separator

\begin{ieltsquestion}
\textbf{Do you prefer to share news with your friends or your parents?}
\end{ieltsquestion}
\textbf{Sample 1: Student Persona (Band 7-8)}

[Nhập câu trả lời vào đây]

\textbf{Sample 2: Working Professional Persona (Band 7-8)}

[Nhập câu trả lời vào đây]

\separator


\vspace{0.8em}
\begin{center}
\renewcommand{\arraystretch}{1.65}
\rowcolors{2}{ForumFreshLight}{white}
\begin{tabularx}{\textwidth}{|V{0.22\textwidth}|C{0.07\textwidth}|I{0.20\textwidth}|Y|V{0.18\textwidth}|}
\hline
\rowcolor{ForumFreshBlue!20}
\textbf{Từ} & \textbf{Loại} & \textbf{IPA} & \textbf{English Meaning} & \textbf{Dịch nghĩa} \\
\hline
\textbf{word} & & /pron/ & & Nghĩa \\ \hline
\end{tabularx}
\end{center}

\newpage

\section{Happy Things}
\begin{ieltsquestion}
\textbf{Is there anything that has made you feel happy lately?}
\end{ieltsquestion}
\textbf{Sample 1: Student Persona (Band 7-8)}

[Nhập câu trả lời vào đây]

\textbf{Sample 2: Working Professional Persona (Band 7-8)}

[Nhập câu trả lời vào đây]

\separator

\begin{ieltsquestion}
\textbf{What made you happy when you were little?}
\end{ieltsquestion}
\textbf{Sample 1: Student Persona (Band 7-8)}

[Nhập câu trả lời vào đây]

\textbf{Sample 2: Working Professional Persona (Band 7-8)}

[Nhập câu trả lời vào đây]

\separator

\begin{ieltsquestion}
\textbf{What do you think will make you feel happy in the future?}
\end{ieltsquestion}
\textbf{Sample 1: Student Persona (Band 7-8)}

[Nhập câu trả lời vào đây]

\textbf{Sample 2: Working Professional Persona (Band 7-8)}

[Nhập câu trả lời vào đây]

\separator

\begin{ieltsquestion}
\textbf{When do you feel happy at work? Why?}
\end{ieltsquestion}
\textbf{Sample 1: Student Persona (Band 7-8)}

[Nhập câu trả lời vào đây]

\textbf{Sample 2: Working Professional Persona (Band 7-8)}

[Nhập câu trả lời vào đây]

\separator

\begin{ieltsquestion}
\textbf{Do you feel happy when buying new things?}
\end{ieltsquestion}
\textbf{Sample 1: Student Persona (Band 7-8)}

[Nhập câu trả lời vào đây]

\textbf{Sample 2: Working Professional Persona (Band 7-8)}

[Nhập câu trả lời vào đây]

\separator

\begin{ieltsquestion}
\textbf{Do you think people are happy when buying new things?}
\end{ieltsquestion}
\textbf{Sample 1: Student Persona (Band 7-8)}

[Nhập câu trả lời vào đây]

\textbf{Sample 2: Working Professional Persona (Band 7-8)}

[Nhập câu trả lời vào đây]

\separator

\begin{ieltsquestion}
\textbf{Do you prefer to watch happy or sad movies?}
\end{ieltsquestion}
\textbf{Sample 1: Student Persona (Band 7-8)}

[Nhập câu trả lời vào đây]

\textbf{Sample 2: Working Professional Persona (Band 7-8)}

[Nhập câu trả lời vào đây]

\separator


\vspace{0.8em}
\begin{center}
\renewcommand{\arraystretch}{1.65}
\rowcolors{2}{ForumFreshLight}{white}
\begin{tabularx}{\textwidth}{|V{0.22\textwidth}|C{0.07\textwidth}|I{0.20\textwidth}|Y|V{0.18\textwidth}|}
\hline
\rowcolor{ForumFreshBlue!20}
\textbf{Từ} & \textbf{Loại} & \textbf{IPA} & \textbf{English Meaning} & \textbf{Dịch nghĩa} \\
\hline
\textbf{word} & & /pron/ & & Nghĩa \\ \hline
\end{tabularx}
\end{center}

\newpage

\section{Doing Something Well}
\begin{ieltsquestion}
\textbf{Do you have an experience when you did something well?}
\end{ieltsquestion}
\textbf{Sample 1: Student Persona (Band 7-8)}

[Nhập câu trả lời vào đây]

\textbf{Sample 2: Working Professional Persona (Band 7-8)}

[Nhập câu trả lời vào đây]

\separator

\begin{ieltsquestion}
\textbf{Do you have an experience when your teacher thought you did a good job?}
\end{ieltsquestion}
\textbf{Sample 1: Student Persona (Band 7-8)}

[Nhập câu trả lời vào đây]

\textbf{Sample 2: Working Professional Persona (Band 7-8)}

[Nhập câu trả lời vào đây]

\separator

\begin{ieltsquestion}
\textbf{Do you often tell your friends when they do something well?}
\end{ieltsquestion}
\textbf{Sample 1: Student Persona (Band 7-8)}

[Nhập câu trả lời vào đây]

\textbf{Sample 2: Working Professional Persona (Band 7-8)}

[Nhập câu trả lời vào đây]

\separator


\vspace{0.8em}
\begin{center}
\renewcommand{\arraystretch}{1.65}
\rowcolors{2}{ForumFreshLight}{white}
\begin{tabularx}{\textwidth}{|V{0.22\textwidth}|C{0.07\textwidth}|I{0.20\textwidth}|Y|V{0.18\textwidth}|}
\hline
\rowcolor{ForumFreshBlue!20}
\textbf{Từ} & \textbf{Loại} & \textbf{IPA} & \textbf{English Meaning} & \textbf{Dịch nghĩa} \\
\hline
\textbf{word} & & /pron/ & & Nghĩa \\ \hline
\end{tabularx}
\end{center}

\newpage

\section{Staying with Old People}
\begin{ieltsquestion}
\textbf{Do you enjoy spending time with old people?}
\end{ieltsquestion}
\textbf{Sample 1: Student Persona (Band 7-8)}

[Nhập câu trả lời vào đây]

\textbf{Sample 2: Working Professional Persona (Band 7-8)}

[Nhập câu trả lời vào đây]

\separator

\begin{ieltsquestion}
\textbf{Have you ever worked with old people?}
\end{ieltsquestion}
\textbf{Sample 1: Student Persona (Band 7-8)}

[Nhập câu trả lời vào đây]

\textbf{Sample 2: Working Professional Persona (Band 7-8)}

[Nhập câu trả lời vào đây]

\separator

\begin{ieltsquestion}
\textbf{Are you happy to work with people who are older than you?}
\end{ieltsquestion}
\textbf{Sample 1: Student Persona (Band 7-8)}

[Nhập câu trả lời vào đây]

\textbf{Sample 2: Working Professional Persona (Band 7-8)}

[Nhập câu trả lời vào đây]

\separator

\begin{ieltsquestion}
\textbf{What are the benefits of being friends with or working with old people?}
\end{ieltsquestion}
\textbf{Sample 1: Student Persona (Band 7-8)}

[Nhập câu trả lời vào đây]

\textbf{Sample 2: Working Professional Persona (Band 7-8)}

[Nhập câu trả lời vào đây]

\separator


\vspace{0.8em}
\begin{center}
\renewcommand{\arraystretch}{1.65}
\rowcolors{2}{ForumFreshLight}{white}
\begin{tabularx}{\textwidth}{|V{0.22\textwidth}|C{0.07\textwidth}|I{0.20\textwidth}|Y|V{0.18\textwidth}|}
\hline
\rowcolor{ForumFreshBlue!20}
\textbf{Từ} & \textbf{Loại} & \textbf{IPA} & \textbf{English Meaning} & \textbf{Dịch nghĩa} \\
\hline
\textbf{word} & & /pron/ & & Nghĩa \\ \hline
\end{tabularx}
\end{center}

\newpage

\section{Life Stages}
\begin{ieltsquestion}
\textbf{Do you enjoy being the age you are now?}
\end{ieltsquestion}
\textbf{Sample 1: Student Persona (Band 7-8)}

[Nhập câu trả lời vào đây]

\textbf{Sample 2: Working Professional Persona (Band 7-8)}

[Nhập câu trả lời vào đây]

\separator

\begin{ieltsquestion}
\textbf{At what age do you think people are the happiest?}
\end{ieltsquestion}
\textbf{Sample 1: Student Persona (Band 7-8)}

[Nhập câu trả lời vào đây]

\textbf{Sample 2: Working Professional Persona (Band 7-8)}

[Nhập câu trả lời vào đây]

\separator

\begin{ieltsquestion}
\textbf{How do people remember each stage of their lives?}
\end{ieltsquestion}
\textbf{Sample 1: Student Persona (Band 7-8)}

[Nhập câu trả lời vào đây]

\textbf{Sample 2: Working Professional Persona (Band 7-8)}

[Nhập câu trả lời vào đây]

\separator

\begin{ieltsquestion}
\textbf{What did you often do with your friends in your childhood?}
\end{ieltsquestion}
\textbf{Sample 1: Student Persona (Band 7-8)}

[Nhập câu trả lời vào đây]

\textbf{Sample 2: Working Professional Persona (Band 7-8)}

[Nhập câu trả lời vào đây]

\separator

\begin{ieltsquestion}
\textbf{Do you have any plans for the next five years?}
\end{ieltsquestion}
\textbf{Sample 1: Student Persona (Band 7-8)}

[Nhập câu trả lời vào đây]

\textbf{Sample 2: Working Professional Persona (Band 7-8)}

[Nhập câu trả lời vào đây]

\separator

\begin{ieltsquestion}
\textbf{What do you think is the most important at the moment?}
\end{ieltsquestion}
\textbf{Sample 1: Student Persona (Band 7-8)}

[Nhập câu trả lời vào đây]

\textbf{Sample 2: Working Professional Persona (Band 7-8)}

[Nhập câu trả lời vào đây]

\separator


\vspace{0.8em}
\begin{center}
\renewcommand{\arraystretch}{1.65}
\rowcolors{2}{ForumFreshLight}{white}
\begin{tabularx}{\textwidth}{|V{0.22\textwidth}|C{0.07\textwidth}|I{0.20\textwidth}|Y|V{0.18\textwidth}|}
\hline
\rowcolor{ForumFreshBlue!20}
\textbf{Từ} & \textbf{Loại} & \textbf{IPA} & \textbf{English Meaning} & \textbf{Dịch nghĩa} \\
\hline
\textbf{word} & & /pron/ & & Nghĩa \\ \hline
\end{tabularx}
\end{center}

\newpage

\section*{PART 1: OBJECTS, SHOPPING \& CONSUMPTION}
\addcontentsline{toc}{section}{PART 1: OBJECTS, SHOPPING \& CONSUMPTION}

\section{Gifts}
\begin{ieltsquestion}
\textbf{What gift have you received recently?}
\end{ieltsquestion}
\textbf{Sample 1: Student Persona (Band 7-8)}

[Nhập câu trả lời vào đây]

\textbf{Sample 2: Working Professional Persona (Band 7-8)}

[Nhập câu trả lời vào đây]

\separator

\begin{ieltsquestion}
\textbf{Have you ever received a great gift?}
\end{ieltsquestion}
\textbf{Sample 1: Student Persona (Band 7-8)}

[Nhập câu trả lời vào đây]

\textbf{Sample 2: Working Professional Persona (Band 7-8)}

[Nhập câu trả lời vào đây]

\separator

\begin{ieltsquestion}
\textbf{Have you ever sent handmade gifts to others?}
\end{ieltsquestion}
\textbf{Sample 1: Student Persona (Band 7-8)}

[Nhập câu trả lời vào đây]

\textbf{Sample 2: Working Professional Persona (Band 7-8)}

[Nhập câu trả lời vào đây]

\separator

\begin{ieltsquestion}
\textbf{What do you consider when choosing a gift?}
\end{ieltsquestion}
\textbf{Sample 1: Student Persona (Band 7-8)}

[Nhập câu trả lời vào đây]

\textbf{Sample 2: Working Professional Persona (Band 7-8)}

[Nhập câu trả lời vào đây]

\separator

\begin{ieltsquestion}
\textbf{Do you think you are good at choosing gifts?}
\end{ieltsquestion}
\textbf{Sample 1: Student Persona (Band 7-8)}

[Nhập câu trả lời vào đây]

\textbf{Sample 2: Working Professional Persona (Band 7-8)}

[Nhập câu trả lời vào đây]

\separator


\vspace{0.8em}
\begin{center}
\renewcommand{\arraystretch}{1.65}
\rowcolors{2}{ForumFreshLight}{white}
\begin{tabularx}{\textwidth}{|V{0.22\textwidth}|C{0.07\textwidth}|I{0.20\textwidth}|Y|V{0.18\textwidth}|}
\hline
\rowcolor{ForumFreshBlue!20}
\textbf{Từ} & \textbf{Loại} & \textbf{IPA} & \textbf{English Meaning} & \textbf{Dịch nghĩa} \\
\hline
\textbf{word} & & /pron/ & & Nghĩa \\ \hline
\end{tabularx}
\end{center}

\newpage

\section{Shoes}
\begin{ieltsquestion}
\textbf{Do you like buying shoes? How often?}
\end{ieltsquestion}
\textbf{Sample 1: Student Persona (Band 7-8)}

[Nhập câu trả lời vào đây]

\textbf{Sample 2: Working Professional Persona (Band 7-8)}

[Nhập câu trả lời vào đây]

\separator

\begin{ieltsquestion}
\textbf{Have you ever bought shoes online?}
\end{ieltsquestion}
\textbf{Sample 1: Student Persona (Band 7-8)}

[Nhập câu trả lời vào đây]

\textbf{Sample 2: Working Professional Persona (Band 7-8)}

[Nhập câu trả lời vào đây]

\separator

\begin{ieltsquestion}
\textbf{How much money do you usually spend on shoes?}
\end{ieltsquestion}
\textbf{Sample 1: Student Persona (Band 7-8)}

[Nhập câu trả lời vào đây]

\textbf{Sample 2: Working Professional Persona (Band 7-8)}

[Nhập câu trả lời vào đây]

\separator

\begin{ieltsquestion}
\textbf{Which do you prefer, fashionable shoes or comfortable shoes?}
\end{ieltsquestion}
\textbf{Sample 1: Student Persona (Band 7-8)}

[Nhập câu trả lời vào đây]

\textbf{Sample 2: Working Professional Persona (Band 7-8)}

[Nhập câu trả lời vào đây]

\separator


\vspace{0.8em}
\begin{center}
\renewcommand{\arraystretch}{1.65}
\rowcolors{2}{ForumFreshLight}{white}
\begin{tabularx}{\textwidth}{|V{0.22\textwidth}|C{0.07\textwidth}|I{0.20\textwidth}|Y|V{0.18\textwidth}|}
\hline
\rowcolor{ForumFreshBlue!20}
\textbf{Từ} & \textbf{Loại} & \textbf{IPA} & \textbf{English Meaning} & \textbf{Dịch nghĩa} \\
\hline
\textbf{word} & & /pron/ & & Nghĩa \\ \hline
\end{tabularx}
\end{center}

\newpage

\section{Advertisement}
\begin{ieltsquestion}
\textbf{Do you like advertisements?}
\end{ieltsquestion}
\textbf{Sample 1: Student Persona (Band 7-8)}

[Nhập câu trả lời vào đây]

\textbf{Sample 2: Working Professional Persona (Band 7-8)}

[Nhập câu trả lời vào đây]

\separator

\begin{ieltsquestion}
\textbf{What kind of advertising do you like?}
\end{ieltsquestion}
\textbf{Sample 1: Student Persona (Band 7-8)}

[Nhập câu trả lời vào đây]

\textbf{Sample 2: Working Professional Persona (Band 7-8)}

[Nhập câu trả lời vào đây]

\separator

\begin{ieltsquestion}
\textbf{Do you often see advertisements when you are on your phone or computer?}
\end{ieltsquestion}
\textbf{Sample 1: Student Persona (Band 7-8)}

[Nhập câu trả lời vào đây]

\textbf{Sample 2: Working Professional Persona (Band 7-8)}

[Nhập câu trả lời vào đây]

\separator

\begin{ieltsquestion}
\textbf{Do you see a lot of advertising on trains or other transport?}
\end{ieltsquestion}
\textbf{Sample 1: Student Persona (Band 7-8)}

[Nhập câu trả lời vào đây]

\textbf{Sample 2: Working Professional Persona (Band 7-8)}

[Nhập câu trả lời vào đây]

\separator

\begin{ieltsquestion}
\textbf{Is there an advertisement that made an impression on you when you were a child?}
\end{ieltsquestion}
\textbf{Sample 1: Student Persona (Band 7-8)}

[Nhập câu trả lời vào đây]

\textbf{Sample 2: Working Professional Persona (Band 7-8)}

[Nhập câu trả lời vào đây]

\separator


\vspace{0.8em}
\begin{center}
\renewcommand{\arraystretch}{1.65}
\rowcolors{2}{ForumFreshLight}{white}
\begin{tabularx}{\textwidth}{|V{0.22\textwidth}|C{0.07\textwidth}|I{0.20\textwidth}|Y|V{0.18\textwidth}|}
\hline
\rowcolor{ForumFreshBlue!20}
\textbf{Từ} & \textbf{Loại} & \textbf{IPA} & \textbf{English Meaning} & \textbf{Dịch nghĩa} \\
\hline
\textbf{word} & & /pron/ & & Nghĩa \\ \hline
\end{tabularx}
\end{center}

\newpage

\section{Borrowing / Lending}
\begin{ieltsquestion}
\textbf{Have you ever borrowed money from others?}
\end{ieltsquestion}
\textbf{Sample 1: Student Persona (Band 7-8)}

[Nhập câu trả lời vào đây]

\textbf{Sample 2: Working Professional Persona (Band 7-8)}

[Nhập câu trả lời vào đây]

\separator

\begin{ieltsquestion}
\textbf{Have you borrowed books from others?}
\end{ieltsquestion}
\textbf{Sample 1: Student Persona (Band 7-8)}

[Nhập câu trả lời vào đây]

\textbf{Sample 2: Working Professional Persona (Band 7-8)}

[Nhập câu trả lời vào đây]

\separator

\begin{ieltsquestion}
\textbf{Do you like to lend things to others?}
\end{ieltsquestion}
\textbf{Sample 1: Student Persona (Band 7-8)}

[Nhập câu trả lời vào đây]

\textbf{Sample 2: Working Professional Persona (Band 7-8)}

[Nhập câu trả lời vào đây]

\separator

\begin{ieltsquestion}
\textbf{Do you mind if others borrow money from you?}
\end{ieltsquestion}
\textbf{Sample 1: Student Persona (Band 7-8)}

[Nhập câu trả lời vào đây]

\textbf{Sample 2: Working Professional Persona (Band 7-8)}

[Nhập câu trả lời vào đây]

\separator

\begin{ieltsquestion}
\textbf{How do you feel when people don't return things they borrowed from you?}
\end{ieltsquestion}
\textbf{Sample 1: Student Persona (Band 7-8)}

[Nhập câu trả lời vào đây]

\textbf{Sample 2: Working Professional Persona (Band 7-8)}

[Nhập câu trả lời vào đây]

\separator


\vspace{0.8em}
\begin{center}
\renewcommand{\arraystretch}{1.65}
\rowcolors{2}{ForumFreshLight}{white}
\begin{tabularx}{\textwidth}{|V{0.22\textwidth}|C{0.07\textwidth}|I{0.20\textwidth}|Y|V{0.18\textwidth}|}
\hline
\rowcolor{ForumFreshBlue!20}
\textbf{Từ} & \textbf{Loại} & \textbf{IPA} & \textbf{English Meaning} & \textbf{Dịch nghĩa} \\
\hline
\textbf{word} & & /pron/ & & Nghĩa \\ \hline
\end{tabularx}
\end{center}

\newpage

\section*{PART 1: MEMORY, DREAMS \& IMAGINATION}
\addcontentsline{toc}{section}{PART 1: MEMORY, DREAMS \& IMAGINATION}

\section{Memory}
\begin{ieltsquestion}
\textbf{Are you good at memorising things?}
\end{ieltsquestion}
\textbf{Sample 1: Student Persona (Band 7-8)}

[Nhập câu trả lời vào đây]

\textbf{Sample 2: Working Professional Persona (Band 7-8)}

[Nhập câu trả lời vào đây]

\separator

\begin{ieltsquestion}
\textbf{Have you ever forgotten something important?}
\end{ieltsquestion}
\textbf{Sample 1: Student Persona (Band 7-8)}

[Nhập câu trả lời vào đây]

\textbf{Sample 2: Working Professional Persona (Band 7-8)}

[Nhập câu trả lời vào đây]

\separator

\begin{ieltsquestion}
\textbf{What do you need to remember in your daily life?}
\end{ieltsquestion}
\textbf{Sample 1: Student Persona (Band 7-8)}

[Nhập câu trả lời vào đây]

\textbf{Sample 2: Working Professional Persona (Band 7-8)}

[Nhập câu trả lời vào đây]

\separator

\begin{ieltsquestion}
\textbf{How do you remember important things?}
\end{ieltsquestion}
\textbf{Sample 1: Student Persona (Band 7-8)}

[Nhập câu trả lời vào đây]

\textbf{Sample 2: Working Professional Persona (Band 7-8)}

[Nhập câu trả lời vào đây]

\separator


\vspace{0.8em}
\begin{center}
\renewcommand{\arraystretch}{1.65}
\rowcolors{2}{ForumFreshLight}{white}
\begin{tabularx}{\textwidth}{|V{0.22\textwidth}|C{0.07\textwidth}|I{0.20\textwidth}|Y|V{0.18\textwidth}|}
\hline
\rowcolor{ForumFreshBlue!20}
\textbf{Từ} & \textbf{Loại} & \textbf{IPA} & \textbf{English Meaning} & \textbf{Dịch nghĩa} \\
\hline
\textbf{word} & & /pron/ & & Nghĩa \\ \hline
\end{tabularx}
\end{center}

\newpage

\section{Dreams}
\begin{ieltsquestion}
\textbf{Can you remember the dreams you had?}
\end{ieltsquestion}
\textbf{Sample 1: Student Persona (Band 7-8)}

[Nhập câu trả lời vào đây]

\textbf{Sample 2: Working Professional Persona (Band 7-8)}

[Nhập câu trả lời vào đây]

\separator

\begin{ieltsquestion}
\textbf{Do you share your dreams with others?}
\end{ieltsquestion}
\textbf{Sample 1: Student Persona (Band 7-8)}

[Nhập câu trả lời vào đây]

\textbf{Sample 2: Working Professional Persona (Band 7-8)}

[Nhập câu trả lời vào đây]

\separator

\begin{ieltsquestion}
\textbf{Do you think dreams have special meanings?}
\end{ieltsquestion}
\textbf{Sample 1: Student Persona (Band 7-8)}

[Nhập câu trả lời vào đây]

\textbf{Sample 2: Working Professional Persona (Band 7-8)}

[Nhập câu trả lời vào đây]

\separator

\begin{ieltsquestion}
\textbf{Do you want to make your dreams come true?}
\end{ieltsquestion}
\textbf{Sample 1: Student Persona (Band 7-8)}

[Nhập câu trả lời vào đây]

\textbf{Sample 2: Working Professional Persona (Band 7-8)}

[Nhập câu trả lời vào đây]

\separator


\vspace{0.8em}
\begin{center}
\renewcommand{\arraystretch}{1.65}
\rowcolors{2}{ForumFreshLight}{white}
\begin{tabularx}{\textwidth}{|V{0.22\textwidth}|C{0.07\textwidth}|I{0.20\textwidth}|Y|V{0.18\textwidth}|}
\hline
\rowcolor{ForumFreshBlue!20}
\textbf{Từ} & \textbf{Loại} & \textbf{IPA} & \textbf{English Meaning} & \textbf{Dịch nghĩa} \\
\hline
\textbf{word} & & /pron/ & & Nghĩa \\ \hline
\end{tabularx}
\end{center}

\newpage

\section{Taking Photos}
\begin{ieltsquestion}
\textbf{Do you like taking photos?}
\end{ieltsquestion}
\textbf{Sample 1: Student Persona (Band 7-8)}

[Nhập câu trả lời vào đây]

\textbf{Sample 2: Working Professional Persona (Band 7-8)}

[Nhập câu trả lời vào đây]

\separator

\begin{ieltsquestion}
\textbf{Do you like taking selfies?}
\end{ieltsquestion}
\textbf{Sample 1: Student Persona (Band 7-8)}

[Nhập câu trả lời vào đây]

\textbf{Sample 2: Working Professional Persona (Band 7-8)}

[Nhập câu trả lời vào đây]

\separator

\begin{ieltsquestion}
\textbf{What is your favourite family photo?}
\end{ieltsquestion}
\textbf{Sample 1: Student Persona (Band 7-8)}

[Nhập câu trả lời vào đây]

\textbf{Sample 2: Working Professional Persona (Band 7-8)}

[Nhập câu trả lời vào đây]

\separator

\begin{ieltsquestion}
\textbf{Do you want to improve your photography skills?}
\end{ieltsquestion}
\textbf{Sample 1: Student Persona (Band 7-8)}

[Nhập câu trả lời vào đây]

\textbf{Sample 2: Working Professional Persona (Band 7-8)}

[Nhập câu trả lời vào đây]

\separator

\begin{ieltsquestion}
\textbf{Do you like to print out your photos or just keep them on your phone or camera?}
\end{ieltsquestion}
\textbf{Sample 1: Student Persona (Band 7-8)}

[Nhập câu trả lời vào đây]

\textbf{Sample 2: Working Professional Persona (Band 7-8)}

[Nhập câu trả lời vào đây]

\separator


\vspace{0.8em}
\begin{center}
\renewcommand{\arraystretch}{1.65}
\rowcolors{2}{ForumFreshLight}{white}
\begin{tabularx}{\textwidth}{|V{0.22\textwidth}|C{0.07\textwidth}|I{0.20\textwidth}|Y|V{0.18\textwidth}|}
\hline
\rowcolor{ForumFreshBlue!20}
\textbf{Từ} & \textbf{Loại} & \textbf{IPA} & \textbf{English Meaning} & \textbf{Dịch nghĩa} \\
\hline
\textbf{word} & & /pron/ & & Nghĩa \\ \hline
\end{tabularx}
\end{center}

\newpage

\section*{PART 1: NATURE, FOOD \& SCHOOL LIFE}
\addcontentsline{toc}{section}{PART 1: NATURE, FOOD \& SCHOOL LIFE}

\section{Food}
\begin{ieltsquestion}
\textbf{What is your favourite food?}
\end{ieltsquestion}
\textbf{Sample 1: Student Persona (Band 7-8)}

[Nhập câu trả lời vào đây]

\textbf{Sample 2: Working Professional Persona (Band 7-8)}

[Nhập câu trả lời vào đây]

\separator

\begin{ieltsquestion}
\textbf{What kind of food did you like when you were young?}
\end{ieltsquestion}
\textbf{Sample 1: Student Persona (Band 7-8)}

[Nhập câu trả lời vào đây]

\textbf{Sample 2: Working Professional Persona (Band 7-8)}

[Nhập câu trả lời vào đây]

\separator

\begin{ieltsquestion}
\textbf{Has your favourite food changed since you were a child?}
\end{ieltsquestion}
\textbf{Sample 1: Student Persona (Band 7-8)}

[Nhập câu trả lời vào đây]

\textbf{Sample 2: Working Professional Persona (Band 7-8)}

[Nhập câu trả lời vào đây]

\separator

\begin{ieltsquestion}
\textbf{Do you eat different foods at different times of the year?}
\end{ieltsquestion}
\textbf{Sample 1: Student Persona (Band 7-8)}

[Nhập câu trả lời vào đây]

\textbf{Sample 2: Working Professional Persona (Band 7-8)}

[Nhập câu trả lời vào đây]

\separator


\vspace{0.8em}
\begin{center}
\renewcommand{\arraystretch}{1.65}
\rowcolors{2}{ForumFreshLight}{white}
\begin{tabularx}{\textwidth}{|V{0.22\textwidth}|C{0.07\textwidth}|I{0.20\textwidth}|Y|V{0.18\textwidth}|}
\hline
\rowcolor{ForumFreshBlue!20}
\textbf{Từ} & \textbf{Loại} & \textbf{IPA} & \textbf{English Meaning} & \textbf{Dịch nghĩa} \\
\hline
\textbf{word} & & /pron/ & & Nghĩa \\ \hline
\end{tabularx}
\end{center}

\newpage

\section{Plants}
\begin{ieltsquestion}
\textbf{Do you keep plants at home?}
\end{ieltsquestion}
\textbf{Sample 1: Student Persona (Band 7-8)}

[Nhập câu trả lời vào đây]

\textbf{Sample 2: Working Professional Persona (Band 7-8)}

[Nhập câu trả lời vào đây]

\separator

\begin{ieltsquestion}
\textbf{What plant did you grow when you were young?}
\end{ieltsquestion}
\textbf{Sample 1: Student Persona (Band 7-8)}

[Nhập câu trả lời vào đây]

\textbf{Sample 2: Working Professional Persona (Band 7-8)}

[Nhập câu trả lời vào đây]

\separator

\begin{ieltsquestion}
\textbf{Do you know anything about growing a plant?}
\end{ieltsquestion}
\textbf{Sample 1: Student Persona (Band 7-8)}

[Nhập câu trả lời vào đây]

\textbf{Sample 2: Working Professional Persona (Band 7-8)}

[Nhập câu trả lời vào đây]

\separator

\begin{ieltsquestion}
\textbf{Do Chinese people send plants as gifts?}
\end{ieltsquestion}
\textbf{Sample 1: Student Persona (Band 7-8)}

[Nhập câu trả lời vào đây]

\textbf{Sample 2: Working Professional Persona (Band 7-8)}

[Nhập câu trả lời vào đây]

\separator


\vspace{0.8em}
\begin{center}
\renewcommand{\arraystretch}{1.65}
\rowcolors{2}{ForumFreshLight}{white}
\begin{tabularx}{\textwidth}{|V{0.22\textwidth}|C{0.07\textwidth}|I{0.20\textwidth}|Y|V{0.18\textwidth}|}
\hline
\rowcolor{ForumFreshBlue!20}
\textbf{Từ} & \textbf{Loại} & \textbf{IPA} & \textbf{English Meaning} & \textbf{Dịch nghĩa} \\
\hline
\textbf{word} & & /pron/ & & Nghĩa \\ \hline
\end{tabularx}
\end{center}

\newpage

\section{Growing Vegetables / Fruits}
\begin{ieltsquestion}
\textbf{Are you interested in growing vegetables and fruits?}
\end{ieltsquestion}
\textbf{Sample 1: Student Persona (Band 7-8)}

[Nhập câu trả lời vào đây]

\textbf{Sample 2: Working Professional Persona (Band 7-8)}

[Nhập câu trả lời vào đây]

\separator

\begin{ieltsquestion}
\textbf{Do you think it's easy to grow vegetables?}
\end{ieltsquestion}
\textbf{Sample 1: Student Persona (Band 7-8)}

[Nhập câu trả lời vào đây]

\textbf{Sample 2: Working Professional Persona (Band 7-8)}

[Nhập câu trả lời vào đây]

\separator

\begin{ieltsquestion}
\textbf{Is growing vegetables popular in your country?}
\end{ieltsquestion}
\textbf{Sample 1: Student Persona (Band 7-8)}

[Nhập câu trả lời vào đây]

\textbf{Sample 2: Working Professional Persona (Band 7-8)}

[Nhập câu trả lời vào đây]

\separator

\begin{ieltsquestion}
\textbf{Do many people grow vegetables in your city?}
\end{ieltsquestion}
\textbf{Sample 1: Student Persona (Band 7-8)}

[Nhập câu trả lời vào đây]

\textbf{Sample 2: Working Professional Persona (Band 7-8)}

[Nhập câu trả lời vào đây]

\separator

\begin{ieltsquestion}
\textbf{Should schools teach students how to grow vegetables?}
\end{ieltsquestion}
\textbf{Sample 1: Student Persona (Band 7-8)}

[Nhập câu trả lời vào đây]

\textbf{Sample 2: Working Professional Persona (Band 7-8)}

[Nhập câu trả lời vào đây]

\separator


\vspace{0.8em}
\begin{center}
\renewcommand{\arraystretch}{1.65}
\rowcolors{2}{ForumFreshLight}{white}
\begin{tabularx}{\textwidth}{|V{0.22\textwidth}|C{0.07\textwidth}|I{0.20\textwidth}|Y|V{0.18\textwidth}|}
\hline
\rowcolor{ForumFreshBlue!20}
\textbf{Từ} & \textbf{Loại} & \textbf{IPA} & \textbf{English Meaning} & \textbf{Dịch nghĩa} \\
\hline
\textbf{word} & & /pron/ & & Nghĩa \\ \hline
\end{tabularx}
\end{center}

\newpage

\section{Pets and Animals}
\begin{ieltsquestion}
\textbf{Have you ever had a pet before?}
\end{ieltsquestion}
\textbf{Sample 1: Student Persona (Band 7-8)}

[Nhập câu trả lời vào đây]

\textbf{Sample 2: Working Professional Persona (Band 7-8)}

[Nhập câu trả lời vào đây]

\separator

\begin{ieltsquestion}
\textbf{What's your favourite animal? Why?}
\end{ieltsquestion}
\textbf{Sample 1: Student Persona (Band 7-8)}

[Nhập câu trả lời vào đây]

\textbf{Sample 2: Working Professional Persona (Band 7-8)}

[Nhập câu trả lời vào đây]

\separator

\begin{ieltsquestion}
\textbf{Where do you prefer to keep your pet, indoors or outdoors?}
\end{ieltsquestion}
\textbf{Sample 1: Student Persona (Band 7-8)}

[Nhập câu trả lời vào đây]

\textbf{Sample 2: Working Professional Persona (Band 7-8)}

[Nhập câu trả lời vào đây]

\separator

\begin{ieltsquestion}
\textbf{What is the most popular animal in China?}
\end{ieltsquestion}
\textbf{Sample 1: Student Persona (Band 7-8)}

[Nhập câu trả lời vào đây]

\textbf{Sample 2: Working Professional Persona (Band 7-8)}

[Nhập câu trả lời vào đây]

\separator


\vspace{0.8em}
\begin{center}
\renewcommand{\arraystretch}{1.65}
\rowcolors{2}{ForumFreshLight}{white}
\begin{tabularx}{\textwidth}{|V{0.22\textwidth}|C{0.07\textwidth}|I{0.20\textwidth}|Y|V{0.18\textwidth}|}
\hline
\rowcolor{ForumFreshBlue!20}
\textbf{Từ} & \textbf{Loại} & \textbf{IPA} & \textbf{English Meaning} & \textbf{Dịch nghĩa} \\
\hline
\textbf{word} & & /pron/ & & Nghĩa \\ \hline
\end{tabularx}
\end{center}

\newpage

\section{Rules}
\begin{ieltsquestion}
\textbf{Are there any rules for students at your school?}
\end{ieltsquestion}
\textbf{Sample 1: Student Persona (Band 7-8)}

[Nhập câu trả lời vào đây]

\textbf{Sample 2: Working Professional Persona (Band 7-8)}

[Nhập câu trả lời vào đây]

\separator

\begin{ieltsquestion}
\textbf{Do you prefer to have more or fewer rules at school?}
\end{ieltsquestion}
\textbf{Sample 1: Student Persona (Band 7-8)}

[Nhập câu trả lời vào đây]

\textbf{Sample 2: Working Professional Persona (Band 7-8)}

[Nhập câu trả lời vào đây]

\separator

\begin{ieltsquestion}
\textbf{Do you think students would benefit more from more rules?}
\end{ieltsquestion}
\textbf{Sample 1: Student Persona (Band 7-8)}

[Nhập câu trả lời vào đây]

\textbf{Sample 2: Working Professional Persona (Band 7-8)}

[Nhập câu trả lời vào đây]

\separator

\begin{ieltsquestion}
\textbf{Have you ever had a really strict teacher?}
\end{ieltsquestion}
\textbf{Sample 1: Student Persona (Band 7-8)}

[Nhập câu trả lời vào đây]

\textbf{Sample 2: Working Professional Persona (Band 7-8)}

[Nhập câu trả lời vào đây]

\separator

\begin{ieltsquestion}
\textbf{Have you ever had a really dedicated teacher?}
\end{ieltsquestion}
\textbf{Sample 1: Student Persona (Band 7-8)}

[Nhập câu trả lời vào đây]

\textbf{Sample 2: Working Professional Persona (Band 7-8)}

[Nhập câu trả lời vào đây]

\separator

\begin{ieltsquestion}
\textbf{Would you like to work as a teacher in a rule-free school?}
\end{ieltsquestion}
\textbf{Sample 1: Student Persona (Band 7-8)}

[Nhập câu trả lời vào đây]

\textbf{Sample 2: Working Professional Persona (Band 7-8)}

[Nhập câu trả lời vào đây]

\separator


\vspace{0.8em}
\begin{center}
\renewcommand{\arraystretch}{1.65}
\rowcolors{2}{ForumFreshLight}{white}
\begin{tabularx}{\textwidth}{|V{0.22\textwidth}|C{0.07\textwidth}|I{0.20\textwidth}|Y|V{0.18\textwidth}|}
\hline
\rowcolor{ForumFreshBlue!20}
\textbf{Từ} & \textbf{Loại} & \textbf{IPA} & \textbf{English Meaning} & \textbf{Dịch nghĩa} \\
\hline
\textbf{word} & & /pron/ & & Nghĩa \\ \hline
\end{tabularx}
\end{center}

\newpage

